\chapter{大规模与在线重建}

\paragraph*{本章阅读锚点}
大规模 3DGS 的问题不是把小场景方法简单跑久一点。原始 3DGS 的优势来自显式高斯、tile-based rasterization 和按视角优化的密度控制；一旦场景扩展到街区、航拍城市、数万张图像或长序列视频，瓶颈会同时出现在四个环节：高斯数量超出单卡显存，远近尺度差异导致无用高斯被反复渲染，分块训练造成边界和遮挡不一致，在线图像流又要求模型边看边更新。本章因此按技术路线阅读：城市级重建看如何分块和处理跨视角遮挡；LoD 路线看如何按距离和 footprint 选择高斯；系统路线看如何把参数放到 CPU、SSD 或多 GPU；在线路线看如何把历史高斯压缩、复用和局部更新。

\section{城市场景重建}

城市路线的基础锚点是原始 3DGS 与 Mega-NeRF/VastGaussian 式分治思想：先把巨大场景切成可训练子问题，再设法把子模型融合成一个可实时渲染的场景。它的难点在于，城市图像通常来自航拍、车载或街景多相机系统，视角覆盖不均、天空和远处区域低纹理、动态物体遮挡严重；如果只按空间网格切块，每个块收到的监督并不一定和最终渲染图像一致。下面这些方法的差别，主要体现在“怎么切”“怎么给每个块分配相机”“怎么合并或配准块”“怎么处理航拍到地面的尺度落差”。



% METHOD: 2403.20032
\paragraph{HO-Gaussian（城市场景的混合初始化与方向编码）（2024.03）}
HO-Gaussian~\cite{2403_20032} 面对的是城市自动驾驶场景中 SfM 初始化不稳定的问题。远处建筑、天空、玻璃和低纹理区域常常没有可靠 SfM 点，原始 3DGS 依赖点云初始化时会出现空洞；多相机车载系统还会带来跨相机外观和几何不一致。

它的改变是引入基于网格的体积表示作为混合优化起点，在场景空间中均匀或自适应采样候选点，再通过梯度驱动 densification 生成高斯，减少对 SfM 点云的依赖。外观上，它用 Gaussian Direction Encoding 替代传统球谐函数，以更紧凑的方向特征表达视角相关颜色。针对多相机系统，方法还加入 neural warping 模块，学习相机间几何变换，缓解同一物体在不同相机中的不一致。相对后续 CityGaussian 系列，HO-Gaussian 更像城市 3DGS 的初始化和外观编码改造，重点是让车载多相机场景先稳定起来。
\textbf{实现细节：} 效率与压缩。EfficientGS（2024.04） 通过选择性致密化和冗余剪枝，在保持渲染保真度的同时将模型大小缩小约 10 倍，使 4K+ 航拍图像的大规模场景重建更加高效。具体而言，该方法在标准 3DGS 致密化过程中引入基于渲染贡献的冗余检测，剪除对图像质量贡献低于阈值的高斯基元，并配合球谐系数降阶进一步压缩存储。HO-Gaussian（2024.03） 提出混合优化方法，将基于网格的体表示与 3DGS 管线结合，消除对 SfM 点初始化的依赖，用高斯方向编码替代球谐函数实现视角依赖颜色表示。CityGS-X（2025.03） 放弃传统的合并-分割过程，采用批级多任务渲染实现高效多 GPU 渲染，4 张 4090 即可在 5 小时内训练渲染 5000+ 图像的场景。
\textbf{方法信息：} Baidu 参与了这篇论文。
\textbf{设计动机与边界：} 相对原始 3DGS 以独立高斯和图像重建误差为主的设定，HO-Gaussian 把问题落在以下可实现的改动上：HO-Gaussian针对城市场景的特殊挑战设计。（1）基于网格体积的点密集化：使用基于网格的体积表示替代SfM初始化。通过在网格上定义密度场来初始化和密集化高斯点，无需依赖SfM点云。网格体积允许在场景空间中均匀采样候选点，并通过梯度驱动密集化增强问题区域的渲染质量。（2）Gaussian Direction Encoding（GDE）：替代传统的球谐函数来编码视角相关颜色。GDE将视角方向编码为更紧凑的特征表示，减少了参数数量同时保持了视角相关外观的表达能力。（3）神经变形模块：针对自动驾驶中多相机系统的特点，引入神经变形来增强不同相机视角间物体的一致性。通过学习相机间的几何变换，减少跨相机视角的渲染不一致。（4）整体管线：端到端可微分，支持实时渲染。其设计目的不是单独增加一个模块，而是用上述表示、监督或渲染约束消除原流程在目标任务中的具体失配。提出HO-Gaussian混合优化方法，结合基于网格的体积与3DGS管线，消除对SfM点初始化的依赖。引入Gaussian Direction Encoding替代球谐函数，实现视角相关颜色表示。提出神经变形（neural warping）模块增强多相机系统中跨相机的物体一致性。适用时还应保留论文给出的边界：网格体积表示在大规模场景中内存开销较大；神经变形模块增加了训练复杂度；未充分处理动态物体建模。
% METHOD: 2407.05254
\paragraph{GaussReg（高斯场之间的快速配准）（2024.07）}
GaussReg~\cite{2407_05254} 是本章的系统代表性锚点之一，但它解决的不是单个 3DGS 的训练，而是多个已经重建的 Gaussian Splatting 场景如何快速对齐。大规模场景常被分段采集和分块重建，若每个块都当作普通点云或图像集合重新做定位，速度和尺度一致性都会成为瓶颈。GaussReg 的改变是把 GS 表示本身当作可配准资产：粗阶段从高斯中抽取低透明度过滤后的点云，用 KPConv-FPN 和 GeoTransformer 做超点匹配，再用 ICP 得到初始相似变换；精阶段利用 GS 可渲染的特点，在粗配准后的重叠区域选择相机并渲染图像，用图像引导的 3D 特征把点云特征和 2D 外观特征合起来做局部精配准。

\textbf{锚点细化：} GaussReg 的输入是两个已经训练好的 GS 场景，以及各自保存的训练相机；输出不是新视角图像，而是把场景 B 对齐到场景 A 的相似变换。粗配准阶段先把高斯均值当作点云，但会根据 opacity 过滤掉弱贡献高斯，避免半透明浮点体干扰匹配；然后用 KPConv-FPN 提取点云局部特征，用 GeoTransformer 建立超点级对应，再通过 ICP 或对应求解得到旋转、平移和尺度的初值。精配准阶段利用 GS 比普通点云多出的一个能力：它可以从任意相机快速渲染外观。方法根据粗变换估计重叠区域，选择两边最有信息量的相机视图并渲染局部图像；2D 图像特征与 3D 点云特征融合成 I3D 特征，重新寻找更精细的局部对应，修正粗阶段的误差。这个设计之所以重要，是因为大场景 3DGS 常需要先局部分块训练，再把块拼成统一地图；GaussReg 把“拼接”从普通点云后处理变成了利用高斯几何和高斯渲染双信息的注册问题。

这个设计相对 HLoc 或纯点云配准的关键差异在于，它既利用高斯的显式几何，又利用高斯能快速渲染的外观监督。输入是两个独立 GS 场景，输出是旋转、平移和尺度组成的相似变换；训练和评测还构建了 ScanNet-GSReg 与 GSReg 数据集。项目 README 与论文标题一致，说明其代码围绕 GS registration pipeline 展开。它不直接降低单场景训练显存，却为城市级拼接提供了一个重要拼图：先让局部 3DGS 成为可稳定对齐的模块，再谈更大的统一表示。
\textbf{方法信息：} Huawei 参与了这篇论文；代码：\url{https://github.com/gap-lab-cuhk-sz/gaussreg}；GitHub 关注度较高，约 135 stars。



% METHOD: 2411.00771
\paragraph{CityGaussianV2（面向大场景的几何精度优化）（2024.11）}
CityGaussianV2~\cite{2411_00771} 继承 CityGaussian 的大规模分治框架，但把注意力转到几何精度。大场景中如果直接用 2DGS 或表面化高斯，表面法向和深度监督能改善几何，却容易因为退化高斯过度拉长而导致数量爆炸；同时，传统基于位置梯度的 densification 对模糊区域和欠重建区域没有区别，容易生成一堆低质量高斯。

它的增量包括三部分。第一，加入单目深度监督来加速几何收敛，使大场景表面不只靠 RGB 对齐。第二，设计 elongation filter，过滤过度拉长或退化的高斯，控制 2DGS 扩展到城市级场景时的数量增长。第三，提出 decomposed-gradient-based densification，把原始位置梯度拆开看待，分别处理模糊伪影和欠覆盖区域。并行训练阶段还结合渲染重要性剪枝、SH 度数蒸馏和 vectree 量化来降低存储。相对 CityGaussian 的“能训练能渲染”，CityGaussianV2 更强调大规模几何基准上的表面准确性。
\textbf{实现细节：} CityGaussianV2（2024.11） 将大规模场景重建从"渲染优先"推进到"渲染与几何并重"的新阶段。基线做法。CityGaussian V1 以 3DGS 为基元，虽然渲染质量出色，但 3D 椭球体的体积表示导致表面定位模糊，提取的 mesh 质量远不能满足实际需求。若将 2DGS 直接替换进 CityGaussian 的框架，会遭遇两个致命问题。（1）收敛困难：2DGS 在训练早期比 3DGS 更容易产生模糊重建。（2）高斯爆炸：2D surfel 从远处观察时投影面积极小（小于 1 像素），经抗锯齿低通滤波后协方差被替换为固定值，无法通过有效梯度调整自身的缩放和旋转。在分块微调中，这些退化的小点被分配到每个块的少量视角频繁观察到，梯度快速累积，触发高斯数量的指数级增长，最终导致内存溢出。改了什么。为解决上述问题，CityGaussianV2 提出三项核心改进。（2）拉伸过滤器（Elongation Filter）：在每次致密化前检查每个 surfel 的长短轴比率，排除极端拉伸的退化 surfel，从源头切断"退化小点→高梯度→致密化爆炸"的正反馈循环。（3）精简训练流水线：将球谐阶数从 3 降至 2，并在块级微调中引入基于百分位贡献评分的在线剪枝（trimming），绕过 CityGaussian 的后处理剪枝和蒸馏阶段，大幅节省存储和内存。此外，还引入了基于 Depth-Anything-V2 的深度先验监督作为几何优化的辅助引导。实现原理。分解梯度致密化（DGD）：标准 3DGS 的致密化梯度为位置梯度的范数 $\nabla_{\mathrm{densify}} = \|\nabla \mathcal{L} / \partial \boldsymbol{\mu}_n\|$。$\max$ 函数中包含常数 $1$ 确保梯度下界不为零。拉伸过滤器（Elongation Filter）：对每个 surfel 定义拉伸率： 其中 $s_{n,u}$ 和 $s_{n,v}$ 分别为 surfel 的两个切向缩放分量。$\rho_n$ 接近 1 意味着极端拉伸（一个轴远长于另一个轴）。在每次致密化操作前，检查所有 surfel 的 $\rho_n$，将 $\rho_n$ 高于某阈值的 surfel 排除在克隆和分裂操作之外。
\textbf{方法信息：} 引用很多，Semantic Scholar 约 55 次。



% METHOD: 2412.01745
\paragraph{Horizon-GS（航拍到地面的统一重建）（2024.12）}
Horizon-GS~\cite{2412_01745} 的局部锚点是城市级 3DGS，但它把问题换成航拍和街景的统一表示。航拍图覆盖范围大、能建立全局几何，但地面细节少；街景图细节丰富，却覆盖局部且数量常常远多于航拍图。如果把两类视角混在一起用原始 3DGS densification，来自俯视和地面视角的梯度会互相冲突，导致高斯既想覆盖大范围结构，又想贴合局部立面细节。

方法的改变是粗到细训练和相机分布平衡。第一阶段主要用航拍图像建立全局粗几何，把城市轮廓、道路拓扑和建筑体量先稳定下来，此时 densification 更偏向补全大尺度结构；第二阶段再引入街景图像补充立面、路面和近景细节，让高斯在局部可见区域细化。为避免街景样本数量压倒航拍样本，它在训练批次中控制不同视角族的比例，使俯视约束不会被大量地面图像淹没。Horizon-GS 还构建航拍-街景配准数据集，使这种训练流程有可评估输入。它相对 BlockGaussian 的差异在于，BlockGaussian 关注空间块内的监督一致，Horizon-GS 关注不同采集高度和视角族之间的 densification 冲突。
\textbf{实现细节：} Horizon-GS（2024.12） 系统性地解决航拍和街景图像的统一重建与渲染问题。在大规模城市场景建模中，现有方法要么仅处理航拍视角（如 VastGaussian、CityGaussian），要么仅处理街景视角（如 Hierarchical-3DGS），两者各自优化但无法协同，导致重建结果在视角切换时出现断裂和伪影。Horizon-GS 通过深入分析发现，将两种视角直接联合训练会引发两大核心冲突：（1）致密化冲突——街景视角促使剪枝遮挡区域的高斯基元，而航拍视角则希望在同一区域重新生长基元，两者梯度方向相互干扰；（2）相机分布不均——街景相机数量远多于航拍相机且覆盖 360 度近景，导致优化过程偏向街景细节、航拍视角严重模糊。为解决上述冲突，Horizon-GS 提出由粗到精的两阶段训练策略。第一阶段仅从航拍图像累积梯度进行致密化，建立覆盖整个场景的粗糙几何骨架，街景视角仅辅助初始高斯基元的全局对齐。第二阶段冻结第一阶段的 MLP 权重和基元属性以保护骨架结构，转而以街景视角为主精化细节。同时引入相机分布平衡机制——在 $[0,1]$ 区间随机采样，若采样值小于 $1/(R+1)$ 则选择航拍图像，否则选择街景图像（第一阶段 $R\!=\!2$ 优先航拍，第二阶段 $R\!=\!1$ 均衡训练）。在致密化策略上，第二阶段不再使用 Scaffold-GS 的平均屏幕空间梯度，而是筛选同时满足高梯度、高不透明度、大投影半径的神经高斯基元作为新锚点，以捕获稀疏观测区域的细粒度细节。该方法还引入多分辨率 LOD 构建策略，分别计算航拍层级数 $K_a$ 和全局层级数 $K$，在第一阶段仅激活 $K_a$ 层、第二阶段开放全部 $K$ 层，以适应两种视角在细节层次上的显著差异。
\textbf{方法信息：} 代码：\url{https://github.com/internrobotics/horizongs}；GitHub 关注度较高，约 145 stars。



% METHOD: 2503.16177
\paragraph{OccluGaussian（遮挡感知场景分割）（2025.03）}
OccluGaussian~\cite{2503_16177} 的锚点是分区域训练，但它指出大规模场景的分块不应只看空间距离，还要看相机之间是否真的共视。城市街区里建筑遮挡强，相机位置相近并不代表它们看到同一组高斯；若把遮挡关系不同的视角分到一个块里，块内优化会收到互相矛盾的监督。

具体做法先构建属性视图图：节点是训练相机，属性包含位置，边权来自 SfM 特征匹配或共视强度。随后用图卷积平滑视图特征，再做谱聚类，把共视强、遮挡关系相近的相机分到同一区域。每个区域独立优化 3DGS，最后合并完整模型。渲染时，方法还能根据视点所在区域保留更可能可见的高斯，从而做区域化加速。相对 BlockGaussian 的点云密度切分，OccluGaussian 的切分依据是视图图和遮挡结构；它更适合遮挡强、街道峡谷式的城市采集。
\textbf{实现细节：} OccluGaussian（2025.03） 关注到现有分割方法忽略遮挡关系——每个子区域可能包含严重遮挡区域，导致区域内的相机相关性低、对整体重建贡献小。该方法提出遮挡感知场景分割策略，其核心是构建属性视角图（attributed view graph）$G = (V, E, X)$：节点 $V$ 对应训练相机，边权 $a_{ij}$ 为 SfM（运动恢复结构，Structure from Motion）特征匹配点数（反映共可见性），节点特征矩阵 $X$ 包含相机三维坐标的位置编码。分割阶段采用属性图聚类算法。先对特征矩阵执行图卷积平滑（滤波器 $G = (I - \frac{1}{2}L_s)^r$，$L_s$ 为对称归一化拉普拉斯矩阵），再对平滑后的相似度矩阵 $\tilde{S} = \frac{1}{2}(|H| + |H^T|)$（$H = \tilde{X}^T\tilde{X}$）进行谱聚类。聚类数通过自适应精化确定：初始值 $K=10$，然后递归分裂过大簇、忽略过小簇或被完全覆盖的簇，直到各簇相机数收敛至 $[M_c - \epsilon_c M_c,\; M_c + \epsilon_c M_c]$（$\epsilon_c = 0.5$）。边界线由 SVM 线性分类器拟合，因遮挡感知聚类而天然对齐场景物理布局。在训练相机选择方面，OccluGaussian 区分三组相机：基础集（位于区域内的相机）、扩展集（区域外但对区域内有充分可见内容的相机）和边界集（朝向区域但被遮挡的相机）。边界集的引入至关重要——缺少这些相机时，边界附近的高斯会变得过大或拉长，导致浮游体和伪影。合并时裁剪掉中心位于区域外的高斯后拼接。同时提出的区域化渲染技术通过预计算可见性掩码 $M_j = \{m_i^j\}_{i=1}^{N_g}$（$m_i^j = 1$ 当且仅当高斯 $i$ 在区域 $R_j$ 任一训练视角下的渲染贡献超过 0.01），在渲染时仅加载当前视点所在区域的可见高斯子集，将遮挡关系转化为显式剔除。
\textbf{实现补充：} 实现上还要核对输入、输出、参数、特征、初始化、训练损失、正则、渲染或推理接口，以及优化中的剪枝、分裂、采样、编码或解码步骤；这些抓手决定该方法到底改了 3DGS 管线的哪一环。
\textbf{方法信息：} Huawei 参与了这篇论文。
\textbf{设计动机与边界：} 相对原始 3DGS 以独立高斯和图像重建误差为主的设定，OccluGaussian 把问题落在以下可实现的改动上：OccluGaussian包含三个核心组件：(1) 遮挡感知场景分割：构建属性视图图，节点为相机（含位置特征），边权重基于SfM中的特征匹配数量表示共视性。使用属性图聚类算法，先进行图卷积生成平滑特征表示，再进行谱聚类分组，确保每个区域内的相机具有强相关性和高平均贡献度。(2) 分区域优化与合并：各区域独立进行3DGS优化，完成后合并为完整模型。(3) 区域化渲染加速：渲染时仅保留视点所在区域可见的高斯，剔除不可见高斯，显著加速渲染而不影响质量。其设计目的不是单独增加一个模块，而是用上述表示、监督或渲染约束消除原流程在目标任务中的具体失配。提出遮挡感知的场景分割策略，基于相机位置和共视性对训练相机进行图聚类，生成具有更强相关性的子区域，并设计区域化渲染加速技术，显著提升大规模场景的重建质量和渲染速度。适用时还应保留论文给出的边界：图聚类过程需要SfM特征匹配结果；场景分割数量需要预设；对于完全没有视觉重叠的区域可能效果有限。
% METHOD: 2503.16964
\paragraph{DroneSplat（野外无人机图像的鲁棒重建）（2025.03）}
DroneSplat~\cite{2503_16964} 也是航拍场景，但它的锚点更接近野外无人机 SfM 加 3DGS。无人机图像会遇到行人、车辆、树叶等动态干扰物，也常因为飞行轨迹稀疏而缺少足够基线。原始 3DGS 会把动态对象吸收到静态场景里，或者在少视角区域生成不稳定高斯；DroneSplat 的改变是同时加入动态 mask 和多视角几何补强。

具体流程先用 SAMv2 得到对象级掩码，再根据渲染残差做局部和全局两级 masking：局部策略在当前训练阶段识别高残差对象，全局策略跨整个序列统计哪些对象持续不稳定，两者合起来剔除动态干扰。对于视角稀疏区域，方法引入 MVS 深度预测作为几何监督，并用体素引导的优化管理高斯密度，让高斯分裂更多发生在确有几何支持的区域。它相对 Horizon-GS 不处理航拍和地面统一，而是把无人机采集中的动态遮挡和稀疏基线作为核心问题。
\textbf{实现细节：} 特殊场景与数据。LetsGo（2024.04） 开发了配备 IMU、LiDAR 和鱼眼相机的手持扫描仪 Polar，构建了 GarageWorld 大型车库数据集，提出深度正则化和多分辨率 3D 高斯表示实现高效 LOD 渲染。Mode-GS（2024.10） 将单目深度锚定引入地面机器人轨迹数据集，通过残差形式高斯解码器和尺度一致深度损失解决地面场景中高斯漂移问题。DroneSplat（2025.03） 针对野外无人机图像中的动态干扰物和有限视角约束，通过局部-全局分割启发式和统计方法自适应调整掩码阈值，结合多视角立体预测和体素引导优化。SplatCo（2025.05） 通过跨结构协作模块融合全局三平面和局部上下文网格特征，配合跨视角剪枝和结构-视角共学习机制实现细节保持的大规模场景渲染。CityGo（2025.05） 采用混合表示——纹理化代理网格加残差高斯——在移动 GPU 上实现实时渲染，训练速度平均提升 1.4 倍。
\textbf{方法信息：} 代码：\url{https://github.com/bityia/dronesplat}；GitHub 关注度较高，约 143 stars。



% METHOD: 2504.09048
\paragraph{BlockGaussian（内容感知分块与可见性监督）（2025.04）}
BlockGaussian~\cite{2504_09048} 继承 CityGaussian/VastGaussian 的分治思路，但它指出普通空间均匀分块会带来两个具体问题：复杂区域和空旷区域负载不均，块内训练图像的监督又常常包含块外内容，导致边界和浮点体伪影。它的改变是把分块、块优化和块合并三个环节都做成可见性感知流程。

具体做法先在曼哈顿世界假设下把稀疏点云对齐到地面平面，将场景投影到水平平面，再用二叉树按点云密度递归切分。点密集区域会被切得更细，空旷区域则保持大块，从而用内容复杂度平衡每块的计算负载。块优化时，它不只用块内点，而是加入辅助点来对齐真实图像里的监督范围，避免渲染图中块外对象被错误解释为块内高斯。合并时又引入伪视角几何约束，从新视角检查空域浮点体并施加约束或移除。相对简单分块，BlockGaussian 的核心是让训练块和图像监督在可见性上对齐；相对 LoD 方法，它主要解决训练可扩展性和块边界一致性，而不是远近渲染开销。
\textbf{实现细节：} BlockGaussian（2025.04） 在分治框架的三个核心环节（分割、训练、合并）上均提出了新的解决方案，特别是针对"分割不均衡""监督不匹配""空间浮游体"三个关键问题给出了系统性回应。Baseline 与问题。在 BlockGaussian 之前，CityGaussian 等方法在收缩空间上做均匀网格划分，虽然缓解了无界场景的空块问题，但未考虑场景内容复杂度的空间差异。建筑密集区域和空旷道路被等同对待，导致各块的重建难度和训练时间差异巨大。此外，分块训练存在两个固有缺陷：（1）监督不匹配——某块的训练视角只能看到该块高斯渲染的部分图像，与完整真值图像做损失计算时，图像中不属于该块的区域会引入错误的梯度信号；（2）空间浮游体——各块独立优化时，缺乏对其占据空间以外的区域（如天空、远距离建筑之间的大段空气空间）的监督，导致大量浮游体产生，合并后严重降低渲染质量。改了什么。BlockGaussian 提出三项核心改进：（1）基于稀疏点云密度分布的内容感知二叉树分割，取代均匀空间划分；（2）辅助高斯 + 小批量优化策略，解决监督不匹配问题；（3）伪视图几何约束损失，在不引入额外训练视角的前提下对空气空间施加间接深度监督，消除浮游体。实现原理。内容感知分割（Content-Aware Scene Partition）：BlockGaussian 利用 SfM 稀疏点云的密度分布作为场景复杂度的估计器。具体流程为：首先基于 Manhattan 世界假设估计地面法向量并沿 $y$ 轴对齐，将稀疏点云投影至 $x$-$z$ 平面并手动定义重建感兴趣区域（Region of Interest, RoI）；然后以二叉树结构递归地分割场景——给定最大树深 $M$ 和叶节点最大点数阈值 $N_b^t$，当某节点深度 $d < M$ 且包含的点云数量 $N_b > N_b^t$ 时，沿该节点包围盒的最长边一分为二，生成两个子节点；否则标记为叶节点。这种自适应分割使得高密度区域（复杂建筑群）被细分为更多小块，低密度区域（开阔路面）保持为大块，自然实现了负载均衡。
\textbf{设计动机与边界：} 相对原始 3DGS 以独立高斯和图像重建误差为主的设定，BlockGaussian 把问题落在以下可实现的改动上：采用分治范式三阶段：(1) 内容感知场景分割：基于曼哈顿世界假设对齐地面平面，将稀疏点云投影到x-z平面，用二叉树结构递归分割场景，根据点云密度（内容复杂度代理）自适应调整分区粒度，确保各块计算负载均衡；(2) 可见性感知块优化：独立优化每个块时引入辅助点（当前块外的点），对齐真值图像的监督信号，避免边界区域的监督不匹配；(3) 伪视角几何约束：在块合并时对空域浮点体（不应存在的高斯）施加几何约束，通过渲染伪视角检测并移除浮点体。最终无缝整合所有块为统一场景表示。其设计目的不是单独增加一个模块，而是用上述表示、监督或渲染约束消除原流程在目标任务中的具体失配。提出BlockGaussian框架实现高效大规模场景重建：(1) 内容感知场景分割策略根据场景复杂度自适应分区，平衡计算负载；(2) 可见性感知块优化引入辅助点对齐真值监督，解决独立块优化的监督不匹配问题；(3) 伪视角几何约束通过空域浮点体消除块合并时的渲染退化。实现5倍优化加速和多基准平均1.21 dB PSNR提升，在单张24GB VRAM设备上实现大规模场景重建。适用时还应保留论文给出的边界：二叉树分割假设场景近似平面；辅助点引入额外计算开销；块边界区域仍可能有轻微质量损失。
\textbf{方法信息：} 代码：\url{https://github.com/sunshinewyc/blockgaussian}；GitHub 关注度较高，约 159 stars。



% METHOD: 2505.17951
\paragraph{SplatCo（结构-视图协同的大场景渲染）（2025.05）}
SplatCo~\cite{2505_17951} 继承 Scaffold-GS 的锚点式表示，但它认为大规模无界场景需要同时保留全局结构和局部细节。普通 3DGS 容易按训练视图过拟合局部纹理，缺少对道路、建筑群和远近层次的全局感知；如果只做剪枝，又可能把对少数视角重要的细节删掉。

方法引入三个模块。跨结构协作模块把全局三平面特征和局部上下文网格特征结合起来，前者提供粗场景布局，后者补充局部纹理和边界。跨视图剪枝机制根据多视角结构一致性移除过拟合高斯，而不是只按不透明度或单视图贡献剪枝。结构视图共学习模块把结构梯度和视图梯度聚合后再更新高斯属性，使几何和外观优化不被单个视角牵着走。相对 OccluGaussian，后者先改相机分组，SplatCo 则在表示内部引入全局-局部协同特征，并配套了面向大规模 splatting 训练和场景处理的 SplatCore 代码。
\textbf{实现细节：} 特殊场景与数据。LetsGo（2024.04） 开发了配备 IMU、LiDAR 和鱼眼相机的手持扫描仪 Polar，构建了 GarageWorld 大型车库数据集，提出深度正则化和多分辨率 3D 高斯表示实现高效 LOD 渲染。Mode-GS（2024.10） 将单目深度锚定引入地面机器人轨迹数据集，通过残差形式高斯解码器和尺度一致深度损失解决地面场景中高斯漂移问题。DroneSplat（2025.03） 针对野外无人机图像中的动态干扰物和有限视角约束，通过局部-全局分割启发式和统计方法自适应调整掩码阈值，结合多视角立体预测和体素引导优化。SplatCo（2025.05） 通过跨结构协作模块融合全局三平面和局部上下文网格特征，配合跨视角剪枝和结构-视角共学习机制实现细节保持的大规模场景渲染。CityGo（2025.05） 采用混合表示——纹理化代理网格加残差高斯——在移动 GPU 上实现实时渲染，训练速度平均提升 1.4 倍。
\textbf{方法信息：} 代码：\url{https://github.com/silas-han/splatcore}。
\textbf{设计动机与边界：} 相对原始 3DGS 以独立高斯和图像重建误差为主的设定，SplatCo 把问题落在以下可实现的改动上：SplatCo框架基于3DGS，针对大规模场景的三个核心挑战（全局结构一致性、局部细节保真度和跨视图一致性）提出了三个关键组件。1）跨结构协作模块（CSCM）：结合全局三平面表示（捕获粗糙场景布局）与局部上下文网格特征（表示精细细节），通过层次补偿机制实现融合，确保全局空间感知和局部细节保持。2）跨视图剪枝机制（CVPM）：基于结构一致性移除过拟合或不准确的高斯，提高存储效率并防止渲染伪影。该机制包含跨视图一致性损失和基于几何一致性的自适应剪枝策略。3）结构视图共学习模块（SVC）：聚合结构梯度与视图梯度，更稳健地引导高斯几何和外观属性的优化。整体框架在Scaffold-GS基础上构建，采用锚点层级架构。其设计目的不是单独增加一个模块，而是用上述表示、监督或渲染约束消除原流程在目标任务中的具体失配。提出跨结构协作模块（CSCM），融合全局三平面特征与局部上下文网格特征；设计跨视图剪枝机制（CVPM），基于结构一致性移除过拟合高斯；引入结构视图共学习模块（SVC），聚合结构与视图梯度以优化高斯属性。在大规模场景渲染中实现了全局结构感知与局部细节保持的平衡。适用时还应保留论文给出的边界：三平面与网格特征的融合需要额外的内存开销；跨视图剪枝可能在高动态场景中误删有效高斯；对于极端大规模场景（如城市级别），分区策略和子模块协调仍有优化空间。
% METHOD: 2604.11401
\paragraph{GS4City（城市模型先验与语义高斯）（2026.04）}
GS4City~\cite{2604_11401} 把城市级 3DGS 从纯视觉重建推向语义可查询场景。它的锚点不是提高 PSNR 的分块策略，而是利用 CityGML/LoD3 这类结构化城市模型，把建筑、墙面、窗户等层级语义注入高斯。普通 3DGS 即便能渲染城市，也不知道某个高斯属于哪栋建筑或哪个构件；后续编辑、查询和数字孪生应用会因此受限。

具体做法先把 CityGML 模型三角化并对齐到 3DGS 坐标系，把 XML 层级展平成建筑-表面-构件三层语义表。然后通过两遍光线投射生成图像对齐掩码：第一遍得到可见几何掩码，第二遍利用父子关系验证细粒度构件，例如窗户必须属于墙面，从而恢复被遮挡或漏检的立面元素。它再把这些几何掩码与 SAM/DINO 等基础模型预测融合，建立场景一致的实例对应，最后为每个高斯学习紧凑身份编码，并用 2D 身份监督和 3D 空间正则优化。相对 SplatCo 的结构特征，GS4City 的结构来自外部城市模型，目标是语义可靠性和可查询性。
\textbf{实现细节：} GS4City（2026.04） 将 CityGML LoD3 标准的结构化城市模型作为几何与语义双重先验引入语义高斯重建，核心改变包括三方面：（1）从 LoD3 城市模型出发，通过两遍光线投射（two-pass raycasting）生成可靠的图像对齐掩码，并利用 CityGML 父子关系恢复单遍投影容易丢失的精细立面元素（如窗户、门洞）；（2）将城市模型掩码与基础模型（SAM + CLIP）生成的非建筑掩码进行优先级融合，通过跨视角几何重叠与视觉-语言语义验证建立场景一致的实例对应关系；（3）为每个高斯学习一个紧凑的 $D$ 维身份编码（identity code），通过 $\alpha$-混合 渲染后由线性分类器预测实例身份，辅以 3D 空间正则化约束近邻语义一致性。实现原理是一条完整的"城市模型→掩码→融合→身份编码"管线。首先将 CityGML 多边形三角化并对齐到重建坐标系，为每个三角面赋予三级层级标签（建筑级 id\_feat、表面级 id\_surf、部件级 id\_part）。第一遍光线投射获取前表面层级标签；第二遍仅对部件级几何投射，通过深度容差 $\epsilon=0.5$\,m 和父表面 ID 一致性验证恢复被墙面遮挡的窗户/门洞。城市模型掩码覆盖建筑区域但不涵盖全场景，因此用 SAM 生成原始掩码后，通过质量过滤和建筑区域重叠比检测，再用 CLIP 视觉-语言模型以边界余量区分前景与建筑。跨视角关联以预训练高斯场作为几何桥梁，结合几何重叠率与语义特征相似度将逐视角掩码聚合为场景一致的全局实例组。融合时采用图像掩码优先规则（非建筑前景优先），并通过固定 ID 偏移量保持城市模型实例与图像实例不冲突。每个融合实例还关联一个跨视角平均的 CLIP 特征，支持开放词汇查询。
\textbf{技术实现补充：} GS4City（2026.04） 将 CityGML LoD3 标准的结构化城市模型作为几何与语义双重先验引入语义高斯重建，核心改变包括三方面：（1）从 LoD3 城市模型出发，通过两遍光线投射（two-pass raycasting）生成可靠的图像对齐掩码，并利用 CityGML 父子关系恢复单遍投影容易丢失的精细立面元素（如窗户、门洞）；（2）将城市模型掩码与基础模型（SAM + CLIP）生成的非建筑掩码进行优先级融合，通过跨视角几何重叠与视觉-语言语义验证建立场景一致的实例对应关系；（3）为每个高斯学习一个紧凑的 $D$ 维身份编码（identity code），通过 $\alpha$-混合 渲染后由线性分类器预测实例身份，辅以 3D 空间正则化约束近邻语义一致性。实现原理是一条完整的"城市模型→掩码→融合→身份编码"管线。首先将 CityGML 多边形三角化并对齐到重建坐标系，为每个三角面赋予三级层级标签（建筑级 id\_feat、表面级 id\_surf、部件级 id\_part）。第一遍光线投射获取前表面层级标签；第二遍仅对部件级几何投射，通过深度容差 $\epsilon=0.5$\,m 和父表面 ID 一致性验证恢复被墙面遮挡的窗户/门洞。城市模型掩码覆盖建筑区域但不涵盖全场景，因此用 SAM 生成原始掩码后，通过质量过滤和建筑区域重叠比检测，再用 CLIP 视觉-语言模型以边界余量区分前景与建筑。跨视角关联以预训练高斯场作为几何桥梁，结合几何重叠率与语义特征相似度将逐视角掩码聚合为场景一致的全局实例组。融合时采用图像掩码优先规则（非建筑前景优先），并通过固定 ID 偏移量保持城市模型实例与图像实例不冲突。每个融合实例还关联一个跨视角平均的 CLIP 特征，支持开放词汇查询。身份编码阶段冻结高斯的几何与外观参数，仅优化每个高斯的 $D$ 维身份向量，通过 $\alpha$-混合 渲染后由线性分类器映射到实例空间，训练目标为逐像素交叉熵损失加上 3D 空间近邻 KL 散度正则项。
\textbf{方法信息：} 代码：\url{https://github.com/jinyzzz/gs4city}。
\textbf{设计动机与边界：} 相对原始 3DGS 以独立高斯和图像重建误差为主的设定，GS4City 把问题落在以下可实现的改动上：给定预训练3DGS场景、校准图像集和对齐城市模型。核心流程：(1)CityGML预处理：三角化LoD3表面、对齐到重建坐标系、将XML层次展平为语义表。每个三角面继承三层ID链(建筑-表面-构件)；(2)两遍光线投射：第一遍获取初始掩码，第二遍利用父-子关系验证细粒度构件(如窗户属于墙壁)，恢复被遮挡的精细立面元素；(3)多源掩码融合：将几何基础的城市模型掩码与基础模型(SAM/DINO)预测融合，建立场景一致的实例对应；(4)高斯身份学习：每个高斯学习紧凑身份编码，在联合2D身份监督和3D空间正则化下优化。其设计目的不是单独增加一个模块，而是用上述表示、监督或渲染约束消除原流程在目标任务中的具体失配。提出GS4City，首个将CityGML城市模型的层次化语义集成到3DGS的方法。通过两遍光线投射从LoD3模型获取图像对齐掩码，利用父子关系验证和恢复精细立面元素。将几何基础掩码与基础模型预测融合建立场景一致实例对应，为每个高斯学习紧凑身份编码。在TUM2TWIN和Gold Coast数据集上，粗粒度建筑分割IoU提升15.8点，细粒度语义分割mIoU提升14.2点。适用时还应保留论文给出的边界：依赖对齐的LoD3城市模型，城市模型与3DGS场景的精确对齐具有挑战性，对城市模型覆盖不到的区域效果有限。

\section{分层与LOD表示}

LoD 路线的共同锚点是一个观察：大场景里并非所有高斯都应在每个视角以最高细节参与渲染。原始 3DGS 只做视锥和 tile 层面的筛选，远处建筑上的细小高斯仍可能被投影、排序和混合，导致 zoom-out 视角速度不稳定。LoD 方法因此把高斯组织成树、金字塔、块级层次或可见性函数，渲染时根据距离、屏幕 footprint、遮挡和缓存状态选择子集。
% METHOD: 2403.17898
\paragraph{Octree-GS（八叉树锚点的 LoD 表示）（2024.03）}
Octree-GS~\cite{2403_17898} 是本节的表示锚点。它继承 Scaffold-GS 的 anchor 思想，但把 anchor 放入八叉树 LoD 层次中。原始 3DGS 和 Scaffold-GS 在大场景 zoom-out 时会把太多细节原语送入渲染，导致速度随视角尺度大幅波动；Octree-GS 的改变是让每个空间位置在多个体素尺度上拥有 anchor，并按观察距离和局部复杂度选择合适层级。

具体做法先从 SfM 点估计相机到点的距离范围，由此决定八叉树层数。LOD 0 是最粗层，层级越高体素越小、细节越多；渲染某一层时采用累积设计，可以包含从粗到细的高斯。每个 anchor 还带有可学习 LOD bias，用来补偿局部高频区域，例如边缘和复杂纹理处会选择更细层级。训练采用从粗到细的渐进式激活，并在每层做 LOD-aware growing 和 pruning。它相对 Mip-Splatting 不是滤掉 aliasing，而是通过空间层次减少不必要的细节渲染；相对 CityGaussian 的块 LoD，它的层次更细粒度地绑定在 anchor 上。
\textbf{实现细节：} 尽管 Scaffold-GS 通过规则间距特征网格改善了结构对齐，Mip-Splatting 通过 3D 平滑和 2D Mip 滤波缓解了高斯冗余，但这些方法在大规模场景中仍然依赖基于可见性的基元筛选——将视锥体内所有基元纳入渲染，而不考虑其投影尺寸。这导致远处的微小细节也被渲染，造成不必要的计算开销和渲染速度的大幅波动。Octree-GS（2024.03） 正是针对这一核心矛盾，提出用八叉树层次结构组织锚点，实现视点距离自适应的 LOD 选择。核心思想：LOD 结构化锚点。Octree-GS 受八叉树数据结构启发，将高斯基元的管理锚点（anchor）组织为层次化的八叉树结构。每个锚点位于不同粒度体素网格的中心——高 LOD 层级的锚点对应更小的体素尺寸，低 LOD 层级对应更大的体素。关键设计是 LOD 的累积性（cumulative）：渲染 LOD $K$ 时，实际光栅化从 LOD 0 到 LOD $K$ 的所有高斯基元。这意味着远处物体使用粗粒度锚点发出的少量大高斯表示，近处物体叠加细粒度锚点提供的精细细节。每个锚点还附带一个可学习的 LOD 偏置（LOD bias）$\Delta L$，用于补偿局部场景复杂性——例如物体锐利边缘区域需要比纯距离判断更高的细节层次。技术细节。Octree-GS 的技术实现包含四个关键组件。（1）\textit{锚点初始化与八叉树构建}。给定 SfM 稀疏点云 $P$，首先根据相机-点距离范围确定八叉树层数 $K = \lfloor\log_2(d_{\max}/d_{\min})\rfloor + 1$，其中 $d_{\max}$ 和 $d_{\min}$ 分别为剔除离群值后的最大和最小观测距离（阈值 $r_d = 0.999$）。然后以基础体素大小 $\delta$ 构建各层体素网格，LOD $L$ 的体素大小为 $\delta / 2^L$，SfM 点落入哪个体素就在该体素中心初始化一个锚点。（2）\textit{视点自适应锚点选择}。
\textbf{设计动机与边界：} 相对原始 3DGS 以独立高斯和图像重建误差为主的设定，Octree-GS 把问题落在以下可实现的改动上：Octree-GS的核心是将锚点按八叉树层级组织，形成自然的LOD层次结构。方法包括：（1）LOD结构化锚点：受Scaffold-GS启发，在不同大小的均匀体素网格中心放置锚点，LOD层级越高（体素越小）捕获的细节越多。LOD设计是累积式的：LOD K渲染时包含LOD 0到K的所有高斯。（2）八叉树初始化：从稀疏SfM点构建八叉树，层数K由相机到SfM点的距离范围决定（K=log2(d\_max/d\_min)+1）。（3）LOD感知的锚点密集化：基于「生长」和「修剪」操作，根据视角信息在需要更多细节的区域增长锚点。（4）渐进式训练策略：从粗到细逐步激活锚点，先训练粗LOD再逐步加入细LOD。（5）外观嵌入：引入外观嵌入以应对野外场景中的光照变化。框架支持多种高斯类型（显式2D/3D高斯或神经高斯）。为了消除原始 3DGS 在该目标任务中的具体失配，方法把上述表示、监督或渲染约束作为可学习和可执行的改动，而不是单独堆叠模块。提出Octree-GS，引入基于八叉树的LOD（细节层次）结构来组织锚点，支持多分辨率场景表示。模型动态从多分辨率锚点集合中选择合适的LOD层级，确保一致的渲染性能和自适应LOD调整，同时保持高保真渲染结果。适用时还应保留论文给出的边界：八叉树构建依赖SfM点质量；对动态场景的支持有限；多LOD层级增加了存储开销；训练策略的超参数（如LOD层级数K）需要根据场景调整。
\textbf{方法信息：} 引用很多，Semantic Scholar 约 287 次。



% METHOD: 2404.01133
\paragraph{CityGaussian（分治训练与跨尺度 LoD 渲染）（2024.04）}
CityGaussian~\cite{2404_01133} 是城市级 LoD 和分治训练的早期锚点。它面对两个同时存在的瓶颈：训练时单卡放不下整个城市，渲染时跨尺度视角又需要不同细节层级。方法的改变是先用全局粗模型提供场景先验，再按先验自适应划分数据和高斯进行块训练，最后通过压缩生成不同 LoD 的高斯集合。

训练阶段，它把无界区域收缩到有界立方体，在全局先验指导下做均匀网格和自适应数据选择，每个块只保留对该块有用的相机视角，降低内存和训练时间。渲染阶段，它为不同块准备不同精度的 Gaussian 集合，并根据相机位置和视距动态选择块级 LoD，再聚合成最终图像。相对 Hierarchical 3DGS 的 tree cut，CityGaussian 更工程化地把“块训练”和“块级 LoD 选择”结合起来；后续 CityGaussianV2 则在此基础上补几何精度和压缩。
\textbf{实现细节：} 方法解决两大挑战：训练内存瓶颈和跨尺度渲染速度。训练阶段采用分而治之策略：先训练全局粗Gaussian先验提供场景全局感知，再基于此先验自适应划分Gaussian和数据进行分块并行训练。将无界区域收缩到归一化有界立方体进行均匀网格划分，平衡各块工作负载。训练数据选择基于全局先验评估视角重要性，只保留对块训练有用的视角。融合阶段合并各块训练结果。渲染阶段引入细节层次(LoD)策略，通过压缩生成不同精度的Gaussian集合。渲染时根据相机位置和视距动态选择各块的合适细节层次，通过分块选择和聚合实现跨尺度实时渲染。例如2.7平方公里的MatrixCity场景23百万Gaussian可实现实时渲染。 提出CityGaussian，采用分而治之训练方法和细节层次(LoD)策略实现高效大尺度3DGS训练和渲染。通过全局场景先验和自适应训练数据选择实现高效训练和无缝融合，基于压缩生成不同细节层次，通过分块细节层次选择和聚合策略实现跨尺度快速渲染。在大尺度场景上达到SOTA渲染质量和一致的实时渲染。 多块并行训练需要多GPU资源；块间边界的无缝融合仍具挑战；LoD层次的过渡可能出现轻微质量跳变；对COLMAP点云质量敏感。
\textbf{实现补充：} 实现上还要核对输入、输出、参数、特征、初始化、训练损失、正则、渲染或推理接口，以及优化中的剪枝、分裂、采样、编码或解码步骤；这些抓手决定该方法到底改了 3DGS 管线的哪一环。
\textbf{方法信息：} 引用很多，Semantic Scholar 约 272 次。
\textbf{设计动机与边界：} 相对原始 3DGS 以独立高斯和图像重建误差为主的设定，CityGaussian 把问题落在以下可实现的改动上：方法解决两大挑战：训练内存瓶颈和跨尺度渲染速度。训练阶段采用分而治之策略：先训练全局粗Gaussian先验提供场景全局感知，再基于此先验自适应划分Gaussian和数据进行分块并行训练。将无界区域收缩到归一化有界立方体进行均匀网格划分，平衡各块工作负载。训练数据选择基于全局先验评估视角重要性，只保留对块训练有用的视角。融合阶段合并各块训练结果。渲染阶段引入细节层次(LoD)策略，通过压缩生成不同精度的Gaussian集合。渲染时根据相机位置和视距动态选择各块的合适细节层次，通过分块选择和聚合实现跨尺度实时渲染。例如2.7平方公里的MatrixCity场景23百万Gaussian可实现实时渲染。其设计目的不是单独增加一个模块，而是用上述表示、监督或渲染约束消除原流程在目标任务中的具体失配。提出CityGaussian，采用分而治之训练方法和细节层次(LoD)策略实现高效大尺度3DGS训练和渲染。通过全局场景先验和自适应训练数据选择实现高效训练和无缝融合，基于压缩生成不同细节层次，通过分块细节层次选择和聚合策略实现跨尺度快速渲染。在大尺度场景上达到SOTA渲染质量和一致的实时渲染。适用时还应保留论文给出的边界：多块并行训练需要多GPU资源；块间边界的无缝融合仍具挑战；LoD层次的过渡可能出现轻微质量跳变；对COLMAP点云质量敏感。
% METHOD: 2405.16829
\paragraph{PyGS（金字塔高斯与视角动态加权）（2024.05）}
PyGS~\cite{2405_16829} 也解决多尺度问题，但它不用八叉树 anchor，而是构造金字塔式高斯集合。顶层由少量大高斯表示低频结构，底层由更多小高斯表示细节；不同视角对不同层的需求不同，因此不能固定渲染全部层。

方法先用快速网格 NeRF 训练一个粗场景，再通过射线投射采样终止点，按不同频率生成多尺度点云来初始化金字塔各层。随后把高斯按空间聚类成若干簇，并训练一个紧凑加权网络，输入相机位姿，输出每个簇在各金字塔层的权重。近距离视角会激活高频层，远距离视角则更依赖顶层大高斯。项目仓库 PyGS README 表明实现侧重 pyramidal scene representation。它相对 Octree-GS 的差异在于，Octree-GS 用显式空间层次选择 anchor，PyGS 用视角条件网络动态混合金字塔层。
\textbf{实现细节：} Horizon-GS（2024.12） 系统性地解决航拍和街景图像的统一重建与渲染问题。在大规模城市场景建模中，现有方法要么仅处理航拍视角（如 VastGaussian、CityGaussian），要么仅处理街景视角（如 Hierarchical-3DGS），两者各自优化但无法协同，导致重建结果在视角切换时出现断裂和伪影。Horizon-GS 通过深入分析发现，将两种视角直接联合训练会引发两大核心冲突：（1）致密化冲突——街景视角促使剪枝遮挡区域的高斯基元，而航拍视角则希望在同一区域重新生长基元，两者梯度方向相互干扰；（2）相机分布不均——街景相机数量远多于航拍相机且覆盖 360 度近景，导致优化过程偏向街景细节、航拍视角严重模糊。为解决上述冲突，Horizon-GS 提出由粗到精的两阶段训练策略。第一阶段仅从航拍图像累积梯度进行致密化，建立覆盖整个场景的粗糙几何骨架，街景视角仅辅助初始高斯基元的全局对齐。第二阶段冻结第一阶段的 MLP 权重和基元属性以保护骨架结构，转而以街景视角为主精化细节。同时引入相机分布平衡机制——在 $[0,1]$ 区间随机采样，若采样值小于 $1/(R+1)$ 则选择航拍图像，否则选择街景图像（第一阶段 $R\!=\!2$ 优先航拍，第二阶段 $R\!=\!1$ 均衡训练）。在致密化策略上，第二阶段不再使用 Scaffold-GS 的平均屏幕空间梯度，而是筛选同时满足高梯度、高不透明度、大投影半径的神经高斯基元作为新锚点，以捕获稀疏观测区域的细粒度细节。该方法还引入多分辨率 LOD 构建策略，分别计算航拍层级数 $K_a$ 和全局层级数 $K$，在第一阶段仅激活 $K_a$ 层、第二阶段开放全部 $K$ 层，以适应两种视角在细节层次上的显著差异。
\textbf{方法信息：} 代码：\url{https://github.com/divyam10/pygs}。
\textbf{设计动机与边界：} 相对原始 3DGS 以独立高斯和图像重建误差为主的设定，PyGS 把问题落在以下可实现的改动上：PyGS针对大规模场景中3DGS的多尺度物体和视角差异问题提出金字塔结构：1）NeRF初始化：先用快速训练的网格NeRF（如TensoRF或InstantNGP）在约20分钟内训练大规模场景，然后通过射线投射采样估计终止点生成点云，在不同频率下采样产生多尺度点云，分别用于初始化金字塔各层的高斯。这克服了COLMAP在大规模场景中耗时长和远处点重建不完整的问题。2）金字塔高斯结构：顶层为少量大尺度高斯捕获低频结构，后续层包含越来越多的较小高斯捕获高频细节。所有层的高斯共同表示场景。3）聚类与动态加权：将金字塔高斯按空间位置聚类为C个簇。设计紧凑的加权网络，输入相机位姿，为每个簇的每个金字塔层级输出一个权重，动态调节各层对该视角的贡献。近视角下高层细节权重更大，远视角下低层结构权重更大。4）端到端训练：所有层的高斯参数和加权网络联合优化。其设计目的不是单独增加一个模块，而是用上述表示、监督或渲染约束消除原流程在目标任务中的具体失配。提出金字塔3DGS（PyGS），以层次化金字塔结构组织高斯。顶层为少量大高斯，后续层为更多小高斯。通过快速训练的网格NeRF在不同频率采样进行初始化，克服COLMAP在大规模场景中的局限。将金字塔高斯聚类分组，用紧凑加权网络根据相机视角动态确定每簇每层的影响力。在多个大规模数据集上达到SOTA，渲染速度比现有方法快400倍以上。适用时还应保留论文给出的边界：NeRF初始化需要额外训练时间；金字塔层数和聚类数需要预设；加权网络引入少量额外计算。
% METHOD: 2406.12080
\paragraph{Hierarchical 3D Gaussian Representation（面向超大数据集的层次高斯）（2024.06）}
Hierarchical 3D Gaussian Representation~\cite{2406_12080} 是本节最重要的 LoD 锚点之一。它继承原始 3DGS 的显式高斯表示，但把问题从“一个场景训练一组高斯”改成“如何让数万张图像、数公里轨迹形成可实时浏览的层次高斯场”。普通分块训练可以把大场景拆开训练，却没有解决渲染时的细节选择：远处建筑仍可能带来大量叶子高斯，块与块之间也容易在层级切换时出现突变。

它的流程分成三步。第一步是 chunk-based 分治训练：把长轨迹或超大图像集切成若干局部块，每个块独立训练叶子层 3DGS，代码仓库中的 \texttt{generate\_chunks.py}、\texttt{train\_single.py} 和相关 SLURM 脚本也反映了这一批处理流程。第二步是 consolidation：把各块训练出的叶子高斯按局部几何和体积属性合并成树状层次，中间节点不是简单平均占位符，而是可以继续优化的高斯集合，用来在远距离视角代表一片区域。第三步是层次渲染：给定相机后，渲染器选择一组 tree cut，近处展开到细叶子，远处停留在粗节点，并在层间做平滑过渡，避免 LoD 层级突然变化。

这个锚点与 Octree-GS 的差异在于，Octree-GS 是把 anchor 放进八叉树层级，训练时逐层激活；Hierarchical 3DGS 则更偏“先局部训练，再构建可优化层次”，目标是把真实大规模采集变成一个可实时漫游的整体。论文报告的场景包含数万张图像、数公里轨迹和多达数十个 chunks；官方仓库提供了 1500 张图像、2 个 chunks 的 toy example 来说明复现步骤。由于这种层次把训练分治、LoD 渲染和资源受限采集放在同一系统里，后续大场景方法经常把它当作 LoD 路线的参照点。
\textbf{实现细节：} Kerbl et al.（2024.06） 提出的层次 3D 高斯表示是 LOD 领域的开创性工作。该方法将高斯基元组织为层次结构，高层节点通过合并低层高斯得到，在渲染远处内容时使用高层粗略表示以提升效率。训练采用分治策略：将大型场景分割为独立块分别训练，然后合并为层次结构并进一步优化中间节点的视觉质量。层次结构的核心构建方式是：对每个父节点高斯，沿其局部坐标系的主轴方向放置 $K_v$ 个子节点。子节点位置偏移由父节点的旋转矩阵和缩放向量决定，几何尺寸按比例因子 $\gamma < 1$ 逐层缩小，同时继承父节点的朝向、颜色和不透明度，从而形成天然的粗到细层次结构。中间节点的优化是该工作的关键创新——不仅优化叶节点，还对合并产生的内部节点进行独立梯度优化，使其视觉质量远超简单继承父节点参数的朴素策略。训练阶段引入多项正则化技术应对大规模稀疏数据，包括：基于屏幕空间覆盖率的自适应密度控制、高斯尺寸下界约束（防止退化）、以及针对远处高斯的不透明度衰减策略。LOD 选择基于每个高斯的 2D 投影半径与屏幕空间阈值的比较，实现自动的细节层级切换和层级间平滑过渡。该方法展示了使用简单设备采集数万张图像、覆盖数公里轨迹、持续长达一小时的实时渲染结果。
\textbf{锚点细化：} Hierarchical 3D Gaussians：针对标准 3DGS 在数千万高斯场景中每次视角切换需处理全部基元的局限，用``沿主轴方向放置 $K_v$ 个子节点的层次合并结构''替代扁平表示，核心思想是\textbf{高层节点通过合并低层高斯得到、远处用粗略表示提升效率}——不仅优化叶节点，还对合并产生的内部节点进行独立梯度优化，使中间节点视觉质量远超朴素继承父节点参数的方案。
\textbf{方法信息：} 代码：\url{https://github.com/graphdeco-inria/hierarchical-3d-gaussians}；GitHub 关注度很高，约 1443 stars；引用很多，Semantic Scholar 约 330 次。



% METHOD: 2505.23158
\paragraph{LODGE（内存受限设备上的块级 LoD）（2025.05）}
LODGE~\cite{2505_23158} 的锚点是移动或内存受限设备上的大场景渲染。很多 LoD 方法假设全部层次数据可常驻 GPU；实际部署时，瓶颈常常是哪些块、哪些层级需要加载。LODGE 的改变是把 LoD 构建、空间分块和动态加载放在同一管线中。

具体做法先从训练好的 3DGS 生成多个细节层级：每个低细节层通过深度感知 3D 平滑滤波、重要性剪枝和短程微调得到，保留整体结构并去除细小高斯。然后把场景划分为空间块，为每块保存多级 LoD。渲染时根据相机位置和距离只加载相关块的合适层级，并用不透明度混合处理块边界，避免硬切换带来的闪烁。相对 Octree-GS 的 anchor 层次，LODGE 更接近部署系统：它关心 GPU 内存、块加载和帧间稳定。
\textbf{实现细节：} 特殊场景与数据。LetsGo（2024.04） 开发了配备 IMU、LiDAR 和鱼眼相机的手持扫描仪 Polar，构建了 GarageWorld 大型车库数据集，提出深度正则化和多分辨率 3D 高斯表示实现高效 LOD 渲染。Mode-GS（2024.10） 将单目深度锚定引入地面机器人轨迹数据集，通过残差形式高斯解码器和尺度一致深度损失解决地面场景中高斯漂移问题。DroneSplat（2025.03） 针对野外无人机图像中的动态干扰物和有限视角约束，通过局部-全局分割启发式和统计方法自适应调整掩码阈值，结合多视角立体预测和体素引导优化。SplatCo（2025.05） 通过跨结构协作模块融合全局三平面和局部上下文网格特征，配合跨视角剪枝和结构-视角共学习机制实现细节保持的大规模场景渲染。CityGo（2025.05） 采用混合表示——纹理化代理网格加残差高斯——在移动 GPU 上实现实时渲染，训练速度平均提升 1.4 倍。
\textbf{方法信息：} Google / DeepMind 参与了这篇论文。
% METHOD: 2507.01110
\paragraph{A LoD of Gaussians（外存层次与无分区训练）（2025.07）}
A LoD of Gaussians~\cite{2507_01110} 反过来质疑分区本身：分区能降低显存，但会引入边界、相机分配和块合并问题。它的目标是在单块消费级 GPU 上训练和渲染超大场景，同时不把场景切成独立子模型。

方法把所有高斯属性放在 CPU 内存中，GPU 只保存树结构副本和当前视角需要的工作集。核心数据结构是高斯层次与序列点树的混合 HSPT：较大子树被转成更高效的序列点树，渲染时用保守的最小距离函数和根-相机距离并行评估 LoD cut，避免选到过粗的节点。缓存系统记录当前驻留 GPU 的 SPT 和细节层次，尽量减少 RAM 到 VRAM 的传输。密集化也在 CPU 侧向层次添加新叶节点或重生低透明叶节点，再转换为 HSPT。相对 LODGE，它不只是动态加载训练后模型，而是把外存、LoD 和 densification 统一到训练过程中。
\textbf{方法信息：} Huawei 参与了这篇论文。
\textbf{设计动机与边界：} 相对原始 3DGS 以独立高斯和图像重建误差为主的设定，A LoD of Gaussians 把问题落在以下可实现的改动上：框架的核心设计包含三个组件：(1) 外存训练：所有高斯属性存储在CPU RAM，根据LoD层次按需流式传输到GPU。仅树结构的副本保持在VRAM中，较大子树替换为序列点树（SPT）形成层次SPT（HSPT）；(2) 混合数据结构：结合高斯层次和SPT，SPT强制更受约束但更高效的切割条件，使用共享的根-相机距离进行并行评估，通过保守最小距离函数避免渲染过粗的高斯。GPU缓存追踪当前驻留的SPT及其细节层次，最小化RAM-VRAM传输；(3) 密集化策略：在CPU上通过向层次添加新叶节点和重生低不透明度叶节点进行密集化，更新后转换为HSPT传输到GPU。支持从航空俯瞰到地面细节的多尺度无缝重建。其设计目的不是单独增加一个模块，而是用上述表示、监督或渲染约束消除原流程在目标任务中的具体失配。提出无需分区的大规模3DGS训练和渲染框架，将全部场景存储在CPU内存中按需流式传输，利用高斯层次与序列点树（SPT）的混合数据结构实现高效视角依赖LoD选择，结合轻量缓存和视图调度系统实现实时流式渲染，支持从航空到地面的多尺度无缝重建。适用时还应保留论文给出的边界：CPU-GPU数据传输可能成为瓶颈；SPT切割条件的约束性可能限制细节表现
% METHOD: 2511.19202
\paragraph{NVGS（神经可见性遮挡剔除）（2025.11）}
NVGS~\cite{2511_19202} 与传统 LoD 互补：它不主要按距离减少细节，而是学习哪些高斯从某个方向看会被遮挡。大规模组合场景中，高斯数量可以超过一亿；即使 LoD 选对了尺度，被前景完全遮挡的后方高斯仍可能进入 rasterization 队列。

具体流程先从预训练 3DGS 资产中采样可见性监督。方法在 Fibonacci 球面上采样方向和距离，对每个视点渲染并判断每个高斯是否对最终像素颜色有贡献，透射率饱和后仍被遮挡的高斯标为不可见。随后把每个高斯的视角依赖可见性函数蒸馏到共享小型 MLP，输入视角方向和距离，输出是否保留。渲染时先做视锥剔除，再用神经可见性剔除，最后通过实例化光栅化器支持组合资产。公开材料包含 trainer 和 viewer 两个仓库，说明它把离线蒸馏和实时查看拆开实现。相对 FilterGS 的层级过滤，NVGS 的过滤依据是遮挡可见性。
\textbf{方法信息：} 代码：\url{https://github.com/UHasselt-DigitalFutureLab/NVGSTrainer}。
\textbf{设计动机与边界：} 相对原始 3DGS 以独立高斯和图像重建误差为主的设定，NVGS 把问题落在以下可实现的改动上：NVGS包含三个组件：(1)可见性提取——从预训练3DGS资产中提取可见性信息，使用Fibonacci球面采样在距离范围内均匀采样2000个方向，对每个视点渲染并判断每个高斯是否对最终像素颜色有贡献（基于透射率饱和度），同时引入随机偏移增强鲁棒性；(2)MLP蒸馏——将每个高斯的可见性函数（输入：视角方向和距离，输出：可见/不可见）蒸馏到小型共享MLP中，利用Tensor Core高效推理；(3)实例化光栅化器——专为组合场景设计，先进行视锥剔除和遮挡剔除，再实例化高斯，支持多资产组合渲染。关键发现：标准体渲染隐含软遮挡（透射率饱和后后续高斯可丢弃），且随距离增加更多高斯被遮挡。其设计目的不是单独增加一个模块，而是用上述表示、监督或渲染约束消除原流程在目标任务中的具体失配。提出NVGS方法，通过小型共享MLP学习高斯的视角依赖可见性函数，在光栅化前丢弃被遮挡原语。配合新型实例化光栅化器，实现超过1亿高斯的实时渲染，在VRAM使用和图像质量上超越SOTA，并与LOD技术互补。适用时还应保留论文给出的边界：MLP推理引入少量延迟；可见性训练需要预训练3DGS资产；对动态场景需重新训练可见性
% METHOD: 2603.23891
\paragraph{FilterGS（无遍历并行过滤）（2026.03）}
FilterGS~\cite{2603_23891} 从 LoD 渲染器本身入手。层次方法常需要沿树遍历为每个视角选择节点，树遍历在 GPU 上不够并行，还会产生大量冗余高斯-瓦片键值对。FilterGS 的改变是用无遍历的并行过滤替代串行 LoD 查询。

它级联两个滤波器：R\&L 滤波器保留叶节点和 footprint 足够小的节点，Ancestor 滤波器沿祖先关系去除同一分支上的冗余层级，最终让每条分支只保留最合适的节点。所有视锥内高斯投影到 2D 后并行处理，不需要逐节点树遍历。方法还定义 Gaussian-Tile Contribution 度量，衡量高斯-瓦片对的冗余程度，并据此做场景自适应高斯收缩，减少排序和混合负担。相对 Octree-GS 关注层次表示怎么建，FilterGS 关注 LoD 表示建好后怎么更快地筛选和 rasterize。
\textbf{实现细节：} FilterGS 提出并行过滤机制——使用两个互补滤波器高效选择高斯元素而无需树遍历。第一个是 Root \& Leaf (R\&L) 滤波器：对视锥体内所有高斯并行计算 2D 投影协方差的特征值，定义像素半径 $R_{2D} = 3\sqrt{\lambda_{\max}}$，当 $R_{2D}$ 小于阈值 $\theta_R$ 时保留该节点（包括叶节点和投影足够小的中间节点），这一步骤完全并行化在 GPU 上执行。第二个是 Ancestor 滤波器：沿祖先路径剪枝冗余节点，确保每个分支仅保留最合适的一个层级。两级滤波器级联的结果是一组可直接送入标准光栅化器的完整高斯基元，无需任何解码或预处理。同时提出 GTC（Gaussian-Tile Count）度量量化高斯-瓦片键值对冗余，基于此的场景自适应高斯收缩策略通过缩小高斯覆盖范围有效减少冗余键值对，在不损失视觉质量的前提下进一步降低光栅化负载。在多个大规模数据集上达到当时最优渲染速度。
\textbf{技术实现补充：} FilterGS 提出并行过滤机制——使用两个互补滤波器高效选择高斯元素而无需树遍历。第一个是 Root \& Leaf (R\&L) 滤波器：对视锥体内所有高斯并行计算 2D 投影协方差的特征值，定义像素半径 $R_{2D} = 3\sqrt{\lambda_{\max}}$，当 $R_{2D}$ 小于阈值 $\theta_R$ 时保留该节点（包括叶节点和投影足够小的中间节点），这一步骤完全并行化在 GPU 上执行。第二个是 Ancestor 滤波器：沿祖先路径剪枝冗余节点，确保每个分支仅保留最合适的一个层级。两级滤波器级联的结果是一组可直接送入标准光栅化器的完整高斯基元，无需任何解码或预处理。同时提出 GTC（Gaussian-Tile Count）度量量化高斯-瓦片键值对冗余，基于此的场景自适应高斯收缩策略通过缩小高斯覆盖范围有效减少冗余键值对，在不损失视觉质量的前提下进一步降低光栅化负载。在多个大规模数据集上达到当时最优渲染速度。
\textbf{方法信息：} 代码：\url{https://github.com/xenon-w/filtergs}。
\textbf{设计动机与边界：} 相对原始 3DGS 以独立高斯和图像重建误差为主的设定，FilterGS 把问题落在以下可实现的改动上：FilterGS包含两个核心组件：(1)无遍历并行过滤：级联两个互补滤波器——R\&L滤波器保留叶子节点和足迹足够小的节点；Ancestor滤波器遍历祖先路径剪枝冗余节点，最终每个分支仅保留最适合的层级。所有视锥内高斯投影到2D后并行处理。(2)GTC（Gaussian-Tile Contribution）度量：量化高斯-瓦片键值对的冗余度。基于此引入场景自适应高斯收缩策略，有效减少冗余对。方法避免了现有LoD方法的两个瓶颈：串行层级遍历（占渲染时间60%以上）和冗余高斯-瓦片处理。其设计目的不是单独增加一个模块，而是用上述表示、监督或渲染约束消除原流程在目标任务中的具体失配。提出FilterGS，包含无遍历并行过滤机制和GTC冗余度量。两个互补滤波器高效选择高斯元素，无需树遍历。引入场景自适应高斯收缩策略减少冗余的高斯-瓦片对。在多个大规模数据集上达到SOTA渲染速度。适用时还应保留论文给出的边界：滤波阈值需要手动设置；在极端大规模场景下的内存管理仍有优化空间。

\section{核外、卸载与大规模训练系统}

这一类方法不一定改变高斯的数学表示，而是改变训练系统的内存层次。原始 3DGS 默认所有高斯参数和优化器状态放在 GPU；到数千万甚至十亿高斯时，参数、梯度和 Adam 状态会远超显存。系统路线的共同思想是利用可见性稀疏性：每个训练视角只激活一小部分高斯，因此 CPU、SSD 或多 GPU 可以承担冷数据存储，GPU 只处理当前工作集。
% METHOD: 2406.18533
\paragraph{On Scaling Up 3D Gaussian Splatting Training（Grendel 式分布式系统）（2024.06）}
On Scaling Up 3D Gaussian Splatting Training~\cite{2406_18533} 是分布式训练路线的高关注锚点候选。它不改变原始 3DGS 的表示和损失，而是把瓶颈定义为“单 GPU 训练系统无法同时容纳更多高斯、更高分辨率图像和更大 batch”。论文的基本观察是，3DGS 各阶段的并行轴并不相同：tile rasterization 更接近图像空间并行，高斯参数和优化器状态更接近模型空间并行；同时，densification 会让高斯数量和空间分布持续变化，静态切分很快失衡。

Grendel-GS 的设计因此采用混合并行。高斯参数被分散到多 GPU，每个设备只常驻一部分高斯；渲染某个相机 batch 时，系统根据空间局部性判断哪些高斯会影响对应像素分区，再通过稀疏 all-to-all 把必要高斯传给相关设备。这样通信量不再是所有设备互发全量模型，而是随当前视角和可见区域变化。训练循环还支持 batch 多视角，而不是原始实现中一次主要处理单视角；学习率等超参数按 batch size 的平方根规则缩放。仓库文件 \texttt{gaussian\_renderer/distribution\_config.py}、\texttt{workload\_division.py} 和多 GPU 示例脚本说明其实现重心确实在分布式调度和负载划分。

它与 DOGS 的关系需要区分清楚：DOGS 把场景切成重叠块，用 ADMM 让局部模型向全局模型达成共识；Grendel-GS 更像保留一套 3DGS 训练目标，但重写训练系统，使高斯参数分区、稀疏通信、动态负载均衡和 batch 多视角训练协同工作。论文报告在 Rubble 场景中，16 GPU 训练 4040 万高斯可比单 GPU 1120 万高斯获得更高 PSNR。也就是说，Grendel-GS 的锚点价值在于说明大规模 3DGS 不只有“分块成多个模型”这一条路，还可以把单一训练过程扩展为分布式参数系统。
\textbf{锚点细化：} Grendel：针对标准 3DGS 训练被限制在单 GPU 造成严重内存瓶颈（A100 40GB 仅容纳约 1120 万高斯）且现有 DNN 分布式框架假设均衡工作负载无法适配 3DGS 不规则稀疏计算的局限，用``高斯-像素混合并行 + 稀疏全到全通信 + 迭代式负载重均衡''替代 FSDP 类密集通信框架，核心思想是\textbf{3DGS 训练管线天然具备混合并行性}——高斯变换阶段适合高斯级分配、图像渲染阶段适合像素级分配，利用空间局部性使阶段间数据交换通过稀疏通信高效完成，16 GPU 在 Rubble 上支持超过 4000 万高斯。
\textbf{方法信息：} 代码：\url{https://github.com/nyu-systems/Grendel-GS}；GitHub 关注度很高，约 676 stars；引用很多，Semantic Scholar 约 57 次。



% METHOD: 2502.10827
\paragraph{E-3DGS（事件相机的大规模无界渲染）（2025.02）}
E-3DGS~\cite{2502_10827} 是本章中比较特殊的系统应用：它把输入从 RGB 帧换成事件流。事件相机只记录亮度变化，适合高速运动和高动态范围场景，但没有普通帧图像那样的 dense photometric supervision；原始 3DGS 的 SfM 初始化和 RGB 损失都不能直接使用。

方法的改变包括四步。第一，在缺少可靠 SfM 点云时，直接在每个相机视锥内随机初始化大量高斯，覆盖无界场景候选空间。第二，把事件流按自适应时间窗口聚合成事件帧，使监督信号随运动速度变化而调整。第三，加入各向同性正则，防止高斯在事件监督不足处退化成极端扁平形状。第四，联合优化相机位姿和高斯参数，修正事件相机轨迹误差。该方法资料中包含官方 E-3DGS 仓库，另有通用 easy-3dgs 工具仓库；本书只把官方仓库列为方法代码。相对普通大规模 RGB 3DGS，E-3DGS 的增量是把渲染监督改造成事件帧重建，并用位姿精化处理传感器特殊性。
\textbf{实现细节：} 与 Grendel 保持原始 3DGS 算法不变、通过系统级并行化加速不同，DOGS（2024.05） 从优化算法层面修改训练过程——将交替方向乘子法（Alternating Direction Method of Multipliers, ADMM）引入 3DGS 训练。场景被递归分解为 $K$ 个有重叠的空间块（每次沿最长轴二分并扩展边界框构造重叠区域），主节点维护全局 3DGS 模型 $z$，$K$ 个从节点维护各自的局部模型 $x_i^k$。每个从节点的训练目标被修改为原渲染损失加正则项：$\min_{x_i^k} L(x_i^k) + \frac{\rho}{2}\|x_i^k - z^t + u_i^k\|_2^{}2$，其中 $u_i^k$ 是对偶变量。正则项的作用是"拉向"全局共识——每个块的局部高斯不能偏离全局模型太远，偏离程度由惩罚参数 $\rho$ 控制。共识步骤通过平均所有从节点的局部高斯更新全局模型 $z^{t+1} = \frac{1}{K}\sum_{k=1}^{K}(x_i^k)^{t+1}$，再将对偶变量推近 $u_i^{t+1} = u_i^t + (x_i^{t+1} - z^{t+1})$。DOGS 引入自适应惩罚参数（根据原始残差和对偶残差的比例动态调整 $\rho$）和过松弛策略（$\alpha^t \in [1.5, 1.8]$）加速 ADMM 收敛，在 2000 次迭代后固定惩罚参数以保证收敛性。从节点训练完成后被丢弃，推理时仅查询全局模型。值得注意的是，DOGS 的递归分割策略沿场景边界框的最长轴二分、重估计每个子块的紧凑边界框后再扩展以构造重叠区域，确保相邻子块在重叠带内共享高斯，从而在 ADMM 共识步骤中对这些共享高斯施加显式一致性约束。实验还揭示了分块数量的一个有趣规律：8 块配置（8 个从节点并行）在大规模场景上不仅训练速度最快，渲染质量也往往优于 4 块或 2 块配置，这是因为更细粒度的空间分解使各子块的局部优化问题更简单、收敛更稳定。
\textbf{方法信息：} 代码：\url{https://github.com/lanpokn/event-3dgs}。
% METHOD: 2509.15645
\paragraph{GS-Scale（主机卸载训练）（2025.09）}
GS-Scale~\cite{2509_15645} 的锚点是主机卸载。它观察到每次 3DGS 迭代只需要当前视锥和 tile 涉及的高斯，但标准实现会把所有参数和 Adam 状态常驻 GPU。方法把大部分高斯属性与优化器状态放到 CPU 主机内存，仅把当前需要的子集传入 GPU。

为了让卸载不拖慢训练，它提出三项优化。选择性卸载把几何属性保留在 GPU，便于快速视锥裁剪，其余颜色、SH 或优化器状态放在主机。参数转发通过提前更新下一步需要的参数，打破 CPU 优化器更新和 GPU 前向/反向之间的串行依赖。延迟优化器更新避免对零梯度高斯做无意义内存访问。公开 GS-Scale README 与论文标题一致。相对 TideGS，GS-Scale 更像 CPU--GPU 双层内存系统；相对 Grendel，它不要求多 GPU 分布式。
\textbf{实现细节：} GS-Scale（2025.09） 将所有高斯存储在主机内存，每次前向/反向传播仅传输子集到 GPU。因此，将占主体的参数和优化器状态卸载到 CPU 可大幅降低显存需求。（2）参数前向（Parameter Forwarding）——提前更新下一迭代所需的高斯参数并传输到 GPU，打破 CPU 优化器更新与 GPU 前向/反向传播之间的数据依赖，实现两者管线化并行执行。GPU 内存需求降低 3.3--5.6 倍，训练速度与纯 GPU 训练可比。
\textbf{实现补充：} 实现上还要核对输入、输出、参数、特征、初始化、训练损失、正则、渲染或推理接口，以及优化中的剪枝、分裂、采样、编码或解码步骤；这些抓手决定该方法到底改了 3DGS 管线的哪一环。
\textbf{方法信息：} 代码：\url{https://github.com/snu-arc/gs-scale}。
\textbf{设计动机与边界：} 相对原始 3DGS 以独立高斯和图像重建误差为主的设定，GS-Scale 把问题落在以下可实现的改动上：GS-Scale系统的核心思路是利用3DGS训练中每次迭代仅小部分高斯参与前向/反向传播的特性：(1) 主机卸载：所有高斯参数和优化器状态存储在主机（CPU）内存，仅按需传输当前需要的子集到GPU；(2) 选择性卸载（Selective Offloading）：仅几何属性（位置、尺度等）保留在GPU用于快速视锥裁剪，其余属性卸载到主机内存；(3) 参数转发（Parameter Forwarding）：通过预更新必要参数，打破CPU优化器更新和GPU前向/反向传播的依赖，实现流水线化；(4) 延迟优化器更新（Deferred Optimizer Update）：延迟零梯度高斯的不必要内存访问，减少CPU内存带宽压力。关键洞察：视锥裁剪和优化器更新被迫在CPU执行，但通过三项优化最小化性能损失。其设计目的不是单独增加一个模块，而是用上述表示、监督或渲染约束消除原流程在目标任务中的具体失配。提出GS-Scale，首个基于主机卸载的大规模3DGS训练系统，三个系统级优化（选择性卸载、参数转发、延迟更新）实现3.3-5.6倍GPU内存降低，训练速度与无卸载GPU相当，使消费级GPU可训练1800万高斯。适用时还应保留论文给出的边界：主机内存带宽成为潜在瓶颈；CPU计算能力限制视锥裁剪和优化器更新速度；参数转发的预更新策略增加实现复杂度；对主机内存容量有较高要求。
% METHOD: 2510.01767
\paragraph{LoBE-GS（负载均衡的大规模训练）（2025.10）}
LoBE-GS~\cite{2510_01767} 关注的是大规模分块训练中的负载不均衡。普通 KD-tree 或均匀网格切分只看空间范围，可能把高斯密集、相机多、遮挡复杂的区域和简单区域分到同等大小块，结果某些块成为训练瓶颈。

它采用负载均衡 KD 树，按 best-first 策略优先切分负载最大的叶节点，并通过 cutline refinement 优化分割线，使两个子块的最大负载差更小。相机分配上，它先由粗模型生成深度图，再把深度反投影到 3D 空间，通过轻量覆盖测试把相机分给相关块，避免传统块-相机全投影的高开销。训练时只保留当前块可见高斯，并在 densification 中跳过已经充分覆盖的高斯，减少冗余分裂。相对 BlockGaussian，LoBE-GS 更系统地把负载估计、相机分配和选择性 densification 连在一起，目标是端到端训练加速。
\textbf{方法信息：} NVIDIA 参与了这篇论文。
\textbf{设计动机与边界：} 相对原始 3DGS 以独立高斯和图像重建误差为主的设定，LoBE-GS 把问题落在以下可实现的改动上：LoBE-GS重新设计大规模3DGS的训练管线，包含三个核心创新：(1)负载均衡KD树场景分割：构建KD树进行场景分区，采用最优优先(best-first)生长策略优先分割负载最大的叶节点。引入分割线优化(cutline refinement)，通过最小化子节点间最大负载差来优化分割位置，实现真正的负载均衡。(2)基于深度的反向投影(DeBPCam)：从粗模型为每个相机计算一次深度图，反向投影到3D空间后通过轻量覆盖测试分配相机到区块，消除传统M×N投影的昂贵计算，将预处理时间从数小时降至数分钟。(3)可见性裁剪(visibility cropping)和选择性致密化(selective densification)：训练时仅保留当前区块可见的Gaussian，并在致密化时跳过被充分覆盖的Gaussian，减少不必要的计算和内存开销。整体采用粗到细的训练范式，粗阶段为细阶段提供初始化。其设计目的不是单独增加一个模块，而是用上述表示、监督或渲染约束消除原流程在目标任务中的具体失配。提出负载均衡的大规模3DGS框架LoBE-GS，通过负载均衡KD树场景分割、基于深度的反向投影快速相机分配、可见性裁剪和选择性致密化三种技术，在保持重建质量的同时实现高达2倍的端到端训练加速。适用时还应保留论文给出的边界：仍依赖多GPU并行训练；粗到细范式需要两阶段训练；负载均衡基于相机数量估计而非实际渲染时间
% METHOD: 2605.20150
\paragraph{TideGS（十亿级高斯的核外训练）（2026.05）}
TideGS~\cite{2605_20150} 是本节的系统代表性锚点。它把 3DGS 的可见性稀疏性推到极致：单个相机视角只访问极少数高斯，因此 GPU VRAM 可以被当作工作集缓存，而不是完整模型存储。方法的目标是在单张 24GB GPU 上训练超过十亿个 Gaussian primitives。

具体做法包括三层。第一，块虚拟化几何把高斯参数按空间块组织，完整参数表常驻 SSD，CPU DRAM 作为中间缓存，GPU 只 materialize 当前相机批次可见的 resident blocks。代码结构中的 \texttt{storage/}、\texttt{tide\_storage\_adapter.py} 和 \texttt{tiered\_cache\_manager.py} 对应的就是这套分层存储抽象。第二，层次异步流水线把 SSD--CPU--GPU 之间的动作重叠起来：CPU 先做块可见性测试，异步预取当前相机需要的块到 GPU，GPU 完成渲染和反传，随后把脏块异步写回 SSD。第三，轨迹自适应差分流只传输相邻迭代间工作集的增量，而不是反复搬运全量参数。

它相对 GS-Scale 的主机卸载更进一步：GS-Scale 主要利用 CPU 内存缓解 GPU 显存，TideGS 把 SSD 纳入训练内存层次，并把“全量模型常驻显存”的假设改成“每步只让可见工作集驻留显存”。论文报告在 MatrixCity 规模场景中训练 1.1B 高斯；素材包记录的实验设置也强调单视图只访问极小比例的高斯。理解 TideGS 时应把它看成训练系统改造，而不是新的渲染方程：高斯的投影、排序和混合仍沿用 3DGS 家族，真正变化的是参数表在哪里、什么时候被搬入 GPU、反向传播后如何写回。
\textbf{实现细节：} TideGS（2026.05） 是首个将 3DGS 训练扩展到十亿高斯量级的单 GPU 核外框架，通过将 SSD--CPU--GPU 构成三级存储层次，使消费级硬件上的超大规模训练成为可能。基线做法。标准 3DGS 要求所有训练状态（参数、梯度、优化器动量）常驻 GPU 显存，单张 24\,GB GPU 仅能容纳约 1100 万高斯。已有的主机卸载方案（CLM、GS-Scale）将部分状态移至 CPU 内存，但仍需在 GPU 上保留关键几何信息或光栅化缓冲区，实际扩展上限约一亿高斯。超过此阈值后，全局基数排序等光栅化辅助缓冲区本身就会耗尽显存（$>$22\,GB）。多 GPU 分布式方案（Grendel-GS、CityGaussian）虽可聚合显存，但需要昂贵的硬件集群和复杂的工程实现，且引入分治边界处理问题。改了什么。TideGS 将 GPU 显存从"持久参数存储"重新定义为"高带宽工作集缓存"——每次迭代仅在 VRAM 中驻留当前相机批次所需的高斯子集，完整参数表始终保存在 SSD 上。这一范式转换将显存需求从 $O(N)$ 解耦为 $O(|R_t|)$（驻留块数量），使单卡可训练规模不受总高斯数限制，而仅受 SSD 容量和带宽的约束。实现原理。TideGS 由三项协同技术支撑 SSD--CPU--GPU 层次管理：第一，块虚拟化几何。将所有高斯参数打包为 SSD 对齐的空间块（$B=4096$ 个高斯/块，约 944\,KiB/块），按 Morton 曲线排序以保证空间邻近高斯映射到相邻块。采用两级可见性过滤：CPU 端对块包围球执行保守视锥体剔除，确定候选工作集 $K_t$；GPU 端在驻留块内执行精确光栅化，确定最终活跃高斯集 $I_t$。块分配在初始化后固定不变，高斯位置更新仅刷新块包围盒但不迁移跨块，因此保持标准 3DGS 渲染语义不变。第二，层次化异步管线把 SSD 读入、CPU 可见性测试、GPU 渲染反传和脏块写回重叠调度；第三，差分数据流只搬运相邻迭代工作集变化，减少重复传输。
\textbf{方法信息：} 代码：\url{https://github.com/sponge-lab/TideGS}；GitHub 关注度较高，约 150 stars。



\section{分布式与并行训练}

分布式路线和系统卸载路线相邻，但侧重点不同：卸载是让一个训练过程容纳更多参数，分布式是让多个设备或节点共同优化同一个大场景。核心困难在于高斯不是独立射线样本；它们有空间重叠、跨块可见和边界一致性问题。好的分布式 3DGS 不只是把数据切开，还要说明局部模型如何对齐成一个全局模型。


% METHOD: 2405.13943
\paragraph{DOGS（ADMM 共识分布式训练）（2024.05）}
DOGS~\cite{2405_13943} 是本章最清晰的分布式代表性锚点。它继承大场景分块训练，但不满足于推理时查询多个子模型；目标是在训练时分布式，推理时仍得到一个全局一致的 3DGS。它的核心改变是把 ADMM 共识优化引入 3DGS：每个从节点训练一个局部块，主节点维护全局高斯，重叠区域的高斯通过共识项保持一致。

具体做法先沿场景最长轴递归二分，把训练视角和点云分成 $K$ 个相交块，并让相邻块保持重叠。每个从节点维护一个局部 3DGS，使用局部图像的 RGB/SSIM 重建损失训练；但对于共享区域的高斯，它还要加入 ADMM 惩罚项，使局部变量 $x_i$ 靠近主节点维护的全局变量 $z$。每个共识步中，主节点收集各从节点的共享高斯，对齐后求平均或聚合得到新的全局状态，再把 $z$ 广播回各块并更新对偶变量。这样，局部训练不会完全漂向各自块内最优，而是被拉回同一个全局模型。

DOGS 与 VastGaussian/BlockGaussian 的关键差异在优化层：后者更强调如何切块、分配相机和合并结果，DOGS 则把“块之间如何一致”写进训练目标。它与 Grendel-GS 的差异也很清楚：Grendel-GS 主要是参数分区和通信调度，目标是扩展一个训练系统；DOGS 直接改变优化形式，把分布式训练建模成共识问题。推理阶段只保留全局 3DGS，不需要像多子模型方案那样查询多个局部模型。官方 README 把项目定位为 NeurIPS 2024 的 Distributed-Oriented Gaussian Splatting，并给出 6 倍以上训练加速的实现说明。
\textbf{实现细节：} 与 Grendel 保持原始 3DGS 算法不变、通过系统级并行化加速不同，DOGS（2024.05） 从优化算法层面修改训练过程——将交替方向乘子法（Alternating Direction Method of Multipliers, ADMM）引入 3DGS 训练。场景被递归分解为 $K$ 个有重叠的空间块（每次沿最长轴二分并扩展边界框构造重叠区域），主节点维护全局 3DGS 模型 $z$，$K$ 个从节点维护各自的局部模型 $x_i^k$。每个从节点的训练目标被修改为原渲染损失加正则项：$\min_{x_i^k} L(x_i^k) + \frac{\rho}{2}\|x_i^k - z^t + u_i^k\|_2^{}2$，其中 $u_i^k$ 是对偶变量。正则项的作用是"拉向"全局共识——每个块的局部高斯不能偏离全局模型太远，偏离程度由惩罚参数 $\rho$ 控制。共识步骤通过平均所有从节点的局部高斯更新全局模型 $z^{t+1} = \frac{1}{K}\sum_{k=1}^{K}(x_i^k)^{t+1}$，再将对偶变量推近 $u_i^{t+1} = u_i^t + (x_i^{t+1} - z^{t+1})$。DOGS 引入自适应惩罚参数（根据原始残差和对偶残差的比例动态调整 $\rho$）和过松弛策略（$\alpha^t \in [1.5, 1.8]$）加速 ADMM 收敛，在 2000 次迭代后固定惩罚参数以保证收敛性。从节点训练完成后被丢弃，推理时仅查询全局模型。值得注意的是，DOGS 的递归分割策略沿场景边界框的最长轴二分、重估计每个子块的紧凑边界框后再扩展以构造重叠区域，确保相邻子块在重叠带内共享高斯，从而在 ADMM 共识步骤中对这些共享高斯施加显式一致性约束。实验还揭示了分块数量的一个有趣规律：8 块配置（8 个从节点并行）在大规模场景上不仅训练速度最快，渲染质量也往往优于 4 块或 2 块配置，这是因为更细粒度的空间分解使各子块的局部优化问题更简单、收敛更稳定。
\textbf{实现补充：} 实现上还要核对输入、输出、参数、特征、初始化、训练损失、正则、渲染或推理接口，以及优化中的剪枝、分裂、采样、编码或解码步骤；这些抓手决定该方法到底改了 3DGS 管线的哪一环。
\textbf{锚点细化：} DOGS：针对 Grendel 保持原始 3DGS 算法不变仅系统级加速、未从优化算法层面修改训练过程的局限，用``ADMM 共识约束 + 自适应惩罚的分布式优化''替代原始 Adam 优化，核心思想是\textbf{将交替方向乘子法引入 3DGS 训练，正则项拉向全局共识}——主节点维护全局模型 $z$，从节点训练局部模型 $x_i^k$ 并在共识步骤平均，自适应惩罚参数和过松弛策略加速 ADMM 收敛，训练速度提升 6 倍以上。
\textbf{方法信息：} 代码：\url{https://github.com/aibluefisher/dogs}；GitHub 关注度较高，约 280 stars；引用很多，Semantic Scholar 约 79 次。



% METHOD: 2412.17612
\paragraph{CoSurfGS（设备-边缘-云协同表面重建）（2024.12）}
CoSurfGS~\cite{2412_17612} 把分布式 3DGS 放到设备-边缘-云架构中，目标是大规模表面重建而非单纯新视角合成。它继承 PGSR/2DGS 一类表面约束思想，但把训练数据分散在多个设备上，例如无人机或边缘节点；直接上传原始图像既占带宽也涉及隐私。

具体流程是设备端先用本地图像训练局部 surface Gaussian 模型，并通过单视图和多视图几何损失把高斯压向表面。上传前，局部模型压缩模块删除重叠区域冗余高斯，降低通信量。模型聚合方案在边缘或云端通过自知识蒸馏保留各局部区域关键结构，同时对齐相邻共视区域的表面几何。设备-边缘和边缘-云使用同一套压缩与聚合操作。相对 DOGS 的 ADMM 共识，CoSurfGS 更关注联邦式或协同式场景下的通信、隐私和表面几何。
\textbf{实现细节：} CoSurfGS（2024.12） 提出多智能体协作的 3DGS 表面重建框架，采用与 Radiant 相似的设备-边缘-云三层分布式架构，但核心目标从渲染质量转向高质量表面几何恢复。CoSurfGS 的两个关键组件是局部模型压缩（LMC）和模型聚合方案（MAS）。LMC 在设备模型上传前执行剪枝：定义优先级评分筛选缺乏多视角一致性且不透明度低的高斯点，移除局部模型在重叠区域的冗余表示，从而大幅降低传输和聚合的内存开销。MAS 采用自知识蒸馏机制渐进式地将局部表面结构聚合为全局模型——在聚合过程中，先用单视角几何约束（法向量邻域一致性）和多视角几何约束（基于单应矩阵的重投影误差与光度一致性）训练每个局部模型，然后通过蒸馏保持各局部区域的关键几何结构并对齐相邻共视区域的表面。
\textbf{方法信息：} Baidu 参与了这篇论文；代码仓库：\url{https://github.com/zju3dv/cosurfgs}，核验时完整实现暂未确认。
\textbf{设计动机与边界：} 相对原始 3DGS 以独立高斯和图像重建误差为主的设定，CoSurfGS 把问题落在以下可实现的改动上：CoSurfGS采用设备-边缘-云三级架构：(1)设备端训练：每个设备（如无人机）捕获图像训练局部高斯模型，使用PGSR的几何约束（单视图和多视图几何损失）将3D高斯扁平化为表面表示，通过平滑权重策略逐步引入几何约束避免收敛困难；(2)局部模型压缩（LMC）：在上传前压缩局部模型，去除重叠区域冗余高斯点，减少传输量和后续聚合开销；(3)模型聚合方案（MAS）：采用自知识蒸馏机制，在聚合局部模型到全局模型时保持各局部区域的关键结构，并对齐相邻共视区域的表面几何。设备-边缘聚合和边缘-云聚合使用相同的f\_MAS和f\_LMC操作。框架仅上传高斯模型而非原始图像，保护设备隐私。其设计目的不是单独增加一个模块，而是用上述表示、监督或渲染约束消除原流程在目标任务中的具体失配。提出CoSurfGS，首个基于分布式学习的大规模3DGS表面重建框架。采用设备-边缘-云架构，通过局部模型压缩（LMC）和模型聚合方案（MAS）实现高质量大规模表面重建，显著降低GPU内存消耗，保护设备隐私。适用时还应保留论文给出的边界：依赖PGSR作为局部训练基线；分布式架构需可靠的通信基础设施；聚合策略在极度重叠区域可能产生冗余。
% METHOD: 2509.12138
\paragraph{Distributed 3DGS for Isosurface Visualization（科学可视化的 HPC 分布式管线）（2025.09）}
Distributed 3D Gaussian Splatting for High-Resolution Isosurface Visualization~\cite{2509_12138} 把 3DGS 用到科学可视化，而不是相机拍摄的自然场景。输入是大型体数据的等值面，目标是在 HPC 系统上训练可交互渲染的高斯表示。普通单 GPU 3DGS 很难容纳上亿级等值面高斯，也不能直接处理分区体数据。

具体管线先用 ParaView 从体数据提取等值面点云作为初始高斯，再生成结构化轨道相机视角。数据集被分成多个分区，每个节点处理一个分区，并添加 ghost cells 以避免分区边界出现渲染缝隙。每个节点渲染自己的分区图像和背景掩码，掩码用于抑制无关 splats 和白色条纹；训练完成后合并各节点的 splats 得到全局模型。项目仓库 Grendel-GS-Scivis 显示其与分布式 GS 科学可视化管线对应。相对城市分布式方法，它的局部锚点是体数据分区和等值面可视化，而非 SfM 相机图像。
\textbf{实现细节：} Han et al.（2025.09） 将分布式 3DGS 扩展到高性能计算（High-Performance Computing, HPC）科学可视化领域。该工作的完整流水线为：先用 ParaView 从体数据中提取等值面点云作为高斯基元初始位置，生成结构化轨道式合成相机视角集（所有节点使用相同设置保证训练一致性），然后将数据按节点分区并添加幽灵单元（ghost cells），各节点独立渲染训练图像和背景掩码后并行训练高斯模型，最终合并所有节点的 splats 进行全局渲染。幽灵单元通过在分区边界向邻域延伸一层额外的数据来消除缝隙和白色条纹伪影，背景掩码则过滤掉对当前分区无关的背景高斯。实验在 Argonne 的 Polaris 超算上进行，使用三种不同规模的科学数据集（Kingsnake 约 4M 点、Rayleigh-Taylor 约 18.2M 点、Richtmyer-Meshkov 约 106.7M 点），每节点配备 4 块 NVIDIA A100 GPU。该工作发表于 SC 2025 Research Poster，为未来原位（in situ）科学可视化奠定了基础。
\textbf{实现补充：} 实现上还要核对输入、输出、参数、特征、初始化、训练损失、正则、渲染或推理接口，以及优化中的剪枝、分裂、采样、编码或解码步骤；这些抓手决定该方法到底改了 3DGS 管线的哪一环。
\textbf{方法信息：} 代码：\url{https://github.com/mengjiaoh/grendel-gs-scivis}。
\textbf{设计动机与边界：} 相对原始 3DGS 以独立高斯和图像重建误差为主的设定，Distributed 3D Gaussian Splatting for High-Resolution Isosurface Visualization 把问题落在以下可实现的改动上：该分布式3DGS管线专为HPC系统设计，流程包括：(1) 等值面提取：使用ParaView从体积数据集提取等值面点云作为初始高斯基元；(2) 相机设置：生成结构化轨道相机视角集，所有节点使用相同设置确保训练一致性；(3) 数据分区：将数据集分为n个分区，每个分区分配到一个计算节点，添加ghost cells（ghost cells）避免分区边界渲染间隙；(4) 图像渲染和掩码：每个节点渲染其数据分区的图像和背景掩码，掩码防止无关splats和白色条纹；(5) 并行训练：每个节点使用多GPU方法独立训练高斯模型；(6) 全局重建：合并所有节点的splats进行最终渲染。该设计无需集中大型数据集，使分布式训练在HPC系统上可行。其设计目的不是单独增加一个模块，而是用上述表示、监督或渲染约束消除原流程在目标任务中的具体失配。提出首个面向HPC的分布式3DGS管线，支持多节点多GPU并行训练，在Polaris超算上8节点实现3倍加速，同时保持图像质量，为大规模科学数据的原位可视化奠定基础。适用时还应保留论文给出的边界：仅验证了等值面可视化场景；合并策略可能引入边界伪影；通信开销在极大规模下可能成为瓶颈；尚未实现真正的原位可视化集成。
% METHOD: 2605.13794
\paragraph{BlitzGS（城市级快速分布式重建）（2026.05）}
BlitzGS~\cite{2605_13794} 的思路是从三个层面削减分布式训练负载。许多分布式大场景方法按空间分块，但空间分块会产生跨块可见性冗余，同一批高斯可能在多个块中被反复处理。BlitzGS 的改变是采用按高斯索引奇偶等非空间方式分片，减少空间边界导致的重复可见性。

在系统层，它将高斯集合分摊到多 GPU，而不是把相邻空间区域完整复制到多个分区。在统计层，它根据高斯对渲染的贡献做重要性排序，裁剪低贡献高斯，降低训练和渲染负担。在渲染层，它自适应选择 tile，跳过无贡献或不可见区域。公开 README 将其定位为 city-scale Gaussian Splatting at lightning speed。相对 DOGS 的共识优化，BlitzGS 更偏负载管理和工程加速；相对 LoBE-GS 的负载均衡 KD 树，它减少的主要是跨空间块可见性冗余。
\textbf{实现细节：} BlitzGS（2026.05） 在 Grendel 的双层并行基础上进一步引入三层工作负载管理，其核心原则是"城市级优化应避免触碰当前设备、当前视角或当前训练阶段所不需要的高斯"。系统层通过索引奇偶性（index parity）而非空间块分片高斯——每张 GPU 存储约 $|G|/M$ 个高斯，渲染时各 GPU 投影本地高斯后通过单次全到全通信将投影高斯路由到所属 tile 的 GPU 上完成光栅化，从根本上缓解了空间分区导致的跨块可见性冗余问题。视角层通过距离 LOD 门控 $L_v(d, t) = \mathrm{clamp}(\lfloor\log_{f}(d_0/d)\rfloor, 0, L_{\max}(t))$ 和基于每视角冗余掩码的贡献性剔除，从几何和重要性两个独立维度裁剪每帧活跃集合。BlitzGS 在城市级场景上实现数量级加速，训练仅需数十分钟。具体而言，索引奇偶性分片使渲染步骤仅需一次跨 GPU 全到全通信（all-to-all exchange）即可完成投影高斯到 tile 所有者的路由，避免了 Grendel 中多轮稀疏全到全通信的延迟累积。在模型层，计划重要性评分与训练进度耦合——早期训练阶段执行较宽松的截断以保留足够的表达容量，后期逐步收紧以移除冗余高斯，确保在数量级加速的同时不损失渲染质量。
\textbf{技术实现补充：} BlitzGS（2026.05） 在 Grendel 的双层并行基础上进一步引入三层工作负载管理，其核心原则是"城市级优化应避免触碰当前设备、当前视角或当前训练阶段所不需要的高斯"。系统层通过索引奇偶性（index parity）而非空间块分片高斯——每张 GPU 存储约 $|G|/M$ 个高斯，渲染时各 GPU 投影本地高斯后通过单次全到全通信将投影高斯路由到所属 tile 的 GPU 上完成光栅化，从根本上缓解了空间分区导致的跨块可见性冗余问题。视角层通过距离 LOD 门控 $L_v(d, t) = \mathrm{clamp}(\lfloor\log_{f}(d_0/d)\rfloor, 0, L_{\max}(t))$ 和基于每视角冗余掩码的贡献性剔除，从几何和重要性两个独立维度裁剪每帧活跃集合。BlitzGS 在城市级场景上实现数量级加速，训练仅需数十分钟。具体而言，索引奇偶性分片使渲染步骤仅需一次跨 GPU 全到全通信（all-to-all exchange）即可完成投影高斯到 tile 所有者的路由，避免了 Grendel 中多轮稀疏全到全通信的延迟累积。在模型层，计划重要性评分与训练进度耦合——早期训练阶段执行较宽松的截断以保留足够的表达容量，后期逐步收紧以移除冗余高斯，确保在数量级加速的同时不损失渲染质量。
\textbf{方法信息：} 代码：\url{https://github.com/akierraee/blitzgs}。
\textbf{设计动机与边界：} 相对原始 3DGS 以独立高斯和图像重建误差为主的设定，BlitzGS 把问题落在以下可实现的改动上：三个层面的工作负载优化：(1)系统级：按高斯索引奇偶性在GPU间分片，而非空间分块。空间分块导致跨块可见性冗余（同一高斯在多个块可见），奇偶分片消除此冗余；(2)统计级：根据高斯对渲染的贡献排序，裁剪低贡献高斯；(3)渲染级：自适应选择渲染瓦片，跳过不可见区域。分布式训练支持多GPU并行。其设计目的不是单独增加一个模块，而是用上述表示、监督或渲染约束消除原流程在目标任务中的具体失配。提出BlitzGS，三个耦合层面的工作负载管理：(1)系统级：按索引奇偶分片（而非空间分块），减少跨块可见性冗余；(2)统计级：高斯重要性排序裁剪冗余；(3)渲染级：自适应瓦片选择。实现城市级3DGS的快速分布式重建。适用时还应保留论文给出的边界：索引奇偶分片可能影响局部优化质量。

\section{在线与增量重建}

在线路线要求模型随图像流更新。它和离线大规模训练的关键差别在于，不能等所有图像到齐后再全局优化；每来一帧都要决定哪些历史高斯保留、压缩、更新或删除。这里的锚点既包括 feed-forward 3DGS，也包括持续学习中的变化检测。
% METHOD: 2503.06235
\paragraph{StreamGS（无位姿图像流的在线重建）（2025.03）}
StreamGS~\cite{2503_06235} 的输入更困难：图像流没有已知相机位姿。它的锚点是 DUSt3R 这类前馈几何模型和 3DGS 的显式渲染；相比离线 COLMAP 加 3DGS，StreamGS 要在流式过程中同时估计相机、生成高斯并去除冗余。

具体流程从初始双视图开始，用 DUSt3R 预测点图和置信度图，并估计相机相对位姿。随后通过内容自适应精炼改善跨帧像素对应关系，尤其处理 DUSt3R 在域外场景中的点图误差。高斯生成分支分别编码几何和外观，将精炼点图和图像特征解码成高斯参数。前馈自适应密度控制根据帧间对应合并重复高斯，减少约四成冗余，同时保持已建场景。相对 LongSplat，它更强调从无位姿输入启动在线 3DGS；相对 CL-Splats，它处理的是连续建图而非已建场景的局部变化。
\textbf{实现细节：} StreamGS（2025.03） 是首个面向无位姿图像流的在线前馈 3DGS 重建方法，能够在不依赖已知相机位姿和内参的条件下，从连续图像流中渐进式构建 3D 高斯表示。基线做法。在线 3DGS 重建面临三重挑战：其一，传统 SfM 预处理需要完整图像集且计算耗时，无法满足实时流式场景；其二，现有可泛化方法（如 pixelSplat、MVSplat）依赖已知位姿进行立体匹配，且仅支持稀疏视角输入（至多 2--3 帧）；其三，无位姿方法（如 CF-3DGS）将位姿作为可学习参数与高斯联合优化，每帧需要数千次迭代，重建速度极慢。改了什么。实现原理。在初始重建阶段，对相邻帧 $(I^{t-1}, I^t)$ 输入 DUSt3R，获得参考坐标系下的像素级 3D 点图及置信度图，再通过点配准问题求解相对位姿（含缩放因子）。在精化阶段，利用 MASt3R 提取局部 3D 特征，通过最近邻互惠搜索（nearest reciprocal）建立帧间稠密匹配，以此重新估计残差变换并更新点图和位姿（PnP-RANSAC）。在高斯生成阶段，用独立的 2D 图像编码器提取外观特征（因 DUSt3R 冻结），将点图坐标与图像特征拼接后送入轻量解码器预测旋转、缩放、不透明度和球谐系数。在合并阶段，利用已建立的 2D 像素对应关系将 3D 空间的高斯合并转化为高效的 2D 像素级特征变形（warping），再由双层卷积网络（MergeNet）将匹配高斯对聚合为单一高斯——整个合并过程无需渲染损失的梯度反传，运行极快。增量式地将每帧新高斯 $G^t$ 与历史重建 $G^{t-1}$ 合并，最终累积为完整场景表示。设计动机。之所以能做到无位姿前馈重建，关键洞察在于：DUSt3R 作为预训练的像素级 3D 点预测器，已从大规模 3D 数据中习得了强几何先验，可直接从图像对预测点图而无需已知位姿。
\textbf{方法信息：} Microsoft 参与了这篇论文。
% METHOD: 2506.21117
\paragraph{CL-Splats（变化检测驱动的持续学习）（2025.06）}
CL-Splats~\cite{2506_21117} 关注的是场景发生变化后的增量更新。原始 3DGS 如果遇到新时间点数据，通常要么从头重训，要么在全模型上继续优化，容易遗忘旧场景或把局部变化扩散到无关区域。CL-Splats 的改变是先定位变化，再只优化变化区域。

给定上一时刻的 3DGS 和新图像集，方法先用 COLMAP 估计新相机姿态，渲染旧模型到新视角，再用 DINOv2 提取真实图像和渲染图的 patch 特征，通过余弦相似度得到 2D 变化掩码。随后把多视角 2D 掩码通过多数表决提升到 3D 高斯掩码，识别哪些高斯需要更新。局部优化阶段只对变化高斯反传，新增高斯表示新物体，删除过时高斯，同时保持不变区域参数。它还保留历史状态，用于回放和时序分割。相对 LongSplat 的前馈流式融合，CL-Splats 更像可持续维护的场景数据库。
\textbf{方法信息：} Google / DeepMind 参与了这篇论文。
\textbf{设计动机与边界：} 相对原始 3DGS 以独立高斯和图像重建误差为主的设定，CL-Splats 把问题落在以下可实现的改动上：核心流程：给定时间t-1的3DGS重建G\^{}t-1和时间t的新图像集，目标是高效更新G\^{}t-1。三大模块：(1)2D变化检测：使用COLMAP估计相机姿态，通过DINOv2提取渲染图像和真实图像的逐块特征图，计算余弦相似度并阈值化得到2D变化掩码，再膨胀处理。(2)3D变化提升：通过掩码多数投票将2D变化掩码提升为3D高斯掩码，识别需要更新的高斯区域。直接用2D掩码约束光度损失效果差，因为缺乏3D空间感知。(3)局部高斯优化：仅在变化区域内优化高斯参数，保持不变区域。新增高斯处理新增内容，删除过时高斯。支持存储和恢复历史场景状态，实现时序分割。其设计目的不是单独增加一个模块，而是用上述表示、监督或渲染约束消除原流程在目标任务中的具体失配。提出首个3DGS持续学习框架，通过鲁棒变化检测和局部优化实现稀疏观测下的增量场景更新，支持历史状态恢复和时序分割，在重建质量和更新效率上超越SOTA。适用时还应保留论文给出的边界：依赖COLMAP姿态估计；变化检测精度受DINOv2特征质量影响；对缓慢渐进变化的检测可能不够敏感。
% METHOD: 2507.16144
\paragraph{LongSplat（长序列在线泛化 3DGS）（2025.07）}
LongSplat~\cite{2507_16144} 是在线路线的代表性锚点。它继承 DUSt3R/DepthSplat 一类前馈多视图重建思路，但针对长序列提出一个核心问题：如果每帧都预测新的高斯并简单累加，历史高斯会快速膨胀且大量冗余；如果压缩过强，又会丢失已经看到的细节。

方法提出 Gaussian-Image Representation，把历史 3D 高斯通过可微透视投影编码成当前视角对齐的类图像 2D 网格。这个表示同时携带历史几何、外观和身份信息，使网络能在当前帧的图像特征与历史高斯状态之间建立对应。每帧处理时，LongSplat 先自适应压缩冗余历史高斯，再把当前视角预测的新高斯在线融合进去；多视角空间特征来自当前帧及邻近帧，历史特征来自 GIR 的浅层 CNN 编码，二者通过 Transformer 融合后输出更新掩码和压缩权重。NVLabs LongSplat README 对应该方法的官方实现说明，正文按该仓库记录代码信息。相对 StreamGS，LongSplat 更强调长序列中的历史状态压缩和身份感知融合。
\textbf{实现细节：} LongSplat（2025.07） 专为长序列图像设计，核心是流式更新机制——增量整合当前视角观测同时选择性压缩冗余历史高斯。动机。标准在线 3DGS 方法逐帧添加高斯但不做冗余管理，高斯数量随序列长度线性增长，长序列场景中很快耗尽内存并导致推理速度急剧下降。关键创新是高斯-图像表示（GIR），将 3D 高斯参数编码为结构化的类图像 2D 格式。具体而言，GIR 通过可微透视投影将历史高斯渲染到当前视角的对齐 2D 网格中，使历史信息与当前帧在同一空间内对齐。每帧的流式更新执行两个关键操作：自适应压缩（选择性压缩历史高斯消除冗余）和在线集成（融合当前视角高斯与历史状态）。特征提取管线结合多视角空间特征（来自当前帧和 $N$ 个时序相邻帧）与历史特征（GIR 编码后经浅层 CNN 提取），通过 Transformer 注意力机制融合空间和时间信息，预测更新掩码和压缩权重。
\textbf{实现补充：} 实现上还要核对输入、输出、参数、特征、初始化、训练损失、正则、渲染或推理接口，以及优化中的剪枝、分裂、采样、编码或解码步骤；这些抓手决定该方法到底改了 3DGS 管线的哪一环。
\textbf{方法信息：} 代码：\url{https://github.com/nvlabs/longsplat}；GitHub 关注度高，约 794 stars。



\section{多GPU与大规模训练抽象}

多 GPU 抽象关注的是研究代码和系统能力之间的落差。大场景 3DGS 不只是高斯更多，还涉及参数分布、tile 负载、显存访问、跨设备同步和上层 PyTorch 接口一致性。与直接提出新分区策略的方法不同，这类工作把多个 GPU 封装成可用的训练/渲染后端，让开发者仍以接近单设备的方式写模型，而底层自动处理高斯和像素工作的分发。



% METHOD: 2606.11390
\paragraph{A Scalable PyTorch Abstraction for Multi-GPU Gaussian Splatting（聚合设备抽象）（2026.06）}
A Scalable PyTorch Abstraction for Multi-GPU Gaussian Splatting~\cite{2606_11390} 关注开发者接口。很多大规模 3DGS 系统需要手写分布式通信、参数切分和自定义 CUDA 调度，导致研究代码难以迁移。该工作的改变是把多 GPU 暴露成一个聚合 PyTorch 设备，模型代码仍像单设备运行，而底层 splatting operator 自动分发高斯参数和渲染计算。

具体实现利用 CUDA unified memory 和 NVLink 管理跨 GPU 数据访问，在 operator 层处理高斯属性的分布、访问和同步。对于上层训练脚本，参数仍表现为普通张量；对于底层后端，渲染器可以把可见高斯和 tile 工作负载分配到多个 GPU。它相对 Grendel 或 DOGS 的差异在于，它不是提出新的分区算法或共识目标，而是提供一个可扩展抽象，降低城市级和高分辨率 GS 训练的工程门槛。
\textbf{实现细节：} A Scalable PyTorch Abstraction for Multi-GPU Gaussian Splatting。A Scalable PyTorch Abstraction for Multi-GPU Gaussian Splatting（2026.06） 延续 Gaussian splatting 在"多GPU与大规模训练"上的基本设定，主要改变是：多GPU PyTorch backend，通过CUDA unified memory和NVLink在operator层分布高斯参数和splatting算子。3DGS 的显式高斯表示和可微渲染能力保持为基础。其设计动机是：提出把多GPU暴露为聚合PyTorch设备的后端，在operator层分布高斯参数和渲染操作，支持城市级重建和高分辨率训练。具体而言，系统利用CUDA unified memory和NVLink管理跨GPU数据访问，模型代码仍像单设备运行，底层splatting operator自动分发计算和参数。适用边界是：依赖NVLink/统一内存硬件条件；跨设备访问模式可能影响性能稳定性。
\textbf{技术实现补充：} A Scalable PyTorch Abstraction for Multi-GPU Gaussian Splatting。A Scalable PyTorch Abstraction for Multi-GPU Gaussian Splatting（2026.06） 延续 Gaussian splatting 在"多GPU与大规模训练"上的基本设定，主要改变是：多GPU PyTorch backend，通过CUDA unified memory和NVLink在operator层分布高斯参数和splatting算子。3DGS 的显式高斯表示和可微渲染能力保持为基础。其设计动机是：提出把多GPU暴露为聚合PyTorch设备的后端，在operator层分布高斯参数和渲染操作，支持城市级重建和高分辨率训练。具体而言，系统利用CUDA unified memory和NVLink管理跨GPU数据访问，模型代码仍像单设备运行，底层splatting operator自动分发计算和参数。适用边界是：依赖NVLink/统一内存硬件条件；跨设备访问模式可能影响性能稳定性。
\textbf{方法信息：} NVIDIA 参与了这篇论文。
\textbf{设计动机与边界：} 相对原始 3DGS 以独立高斯和图像重建误差为主的设定，A Scalable PyTorch Abstraction for Multi-GPU Gaussian Splatting 把问题落在以下可实现的改动上：系统利用CUDA unified memory和NVLink管理跨GPU数据访问，模型代码仍像单设备运行，底层splatting operator自动分发计算和参数。其设计目的不是单独增加一个模块，而是用上述表示、监督或渲染约束消除原流程在目标任务中的具体失配。提出把多GPU暴露为聚合PyTorch设备的后端，在operator层分布高斯参数和渲染操作，支持城市级重建和高分辨率训练。适用时还应保留论文给出的边界：依赖NVLink/统一内存硬件条件；跨设备访问模式可能影响性能稳定性。

\section{数字孪生与通信建模}

数字孪生应用常把 3DGS 的“可渲染”能力转化为更一般的空间场建模能力。它们未必输出 RGB 新视角，而是利用高斯核的局部连续性、可插值性和几何贴附性来表示无线信道、建筑环境或工程状态。本节的读法应与视觉重建不同：重点不是 PSNR，而是高斯是否为真实空间系统提供了可查询、可更新的局部统计模型。
% METHOD: 2606.18734
\paragraph{Tangent Gaussian Splatting for Channel Prediction（点云辅助无线信道预测）（2026.06）}
Point-Cloud-Assistant Localized Statistical Channel Prediction by Tangent Gaussian Splatting~\cite{2606_18734} 是 3DGS 在通信数字孪生中的边界案例。它不以新视角合成为目标，而是用点云几何和局部高斯表示来预测未测位置的无线信道统计量，例如角功率谱。这里的锚点是传统 ray tracing 或统计信道建模；高斯的作用是提供局部连续、可插值的空间表示。

方法围绕参考信号接收功率和点云几何构建局部 tangent Gaussian 表示，把建筑、墙面或散射体附近的统计信息编码到局部高斯中。查询未测位置时，系统根据点云辅助几何和高斯局部统计预测多径角功率谱，从而服务网络规划或定位。相对本章其他方法，它没有改 3DGS 渲染器来生成 RGB 图像，而是借用 Gaussian Splatting 的局部核表示和可微插值思想；因此适合作为数字孪生应用的简短收录，而不是大规模视觉重建锚点。
具体做法上，Point-Cloud-Assistant Localized Statistical Channel Predicti 将本节问题落实为以下流程：围绕参考信号接收功率和点云几何构建局部高斯表示，预测多径角功率谱。
\textbf{实现细节：} Point-Cloud-Assistant Localized Statistical Channel Prediction by Tangent Gaussian Splatting（2026.06） 延续 Gaussian splatting 在"数字孪生与通信建模"上的基本设定，主要改变是：Tangent Gaussian Splatting辅助的localized statistical channel prediction。3DGS 的显式高斯表示和可微渲染能力保持为基础。其设计动机是：利用点云和高斯局部表示补全未测位置的信道统计信息，服务大规模网络优化。具体而言，方法围绕参考信号接收功率和点云几何构建局部高斯表示，预测多径角功率谱。适用边界是：任务不是视觉3DGS重建；高斯splatting更多作为几何/统计建模工具。
\textbf{实现补充：} 实现上还要核对输入、输出、参数、特征、初始化、训练损失、正则、渲染或推理接口，以及优化中的剪枝、分裂、采样、编码或解码步骤；这些抓手决定该方法到底改了 3DGS 管线的哪一环。
\textbf{方法信息：} Huawei 参与了这篇论文。
\section{结构感知调度}

结构感知调度回到大规模训练中的时间安排问题。原始 3DGS 依赖图像残差和梯度触发 densification，但在稀疏观测区域，高频图像监督可能先于可靠几何到来，导致高斯数量膨胀或漂浮。本节方法不是换一个新表示，而是调度监督频率、图像采样和密度增长，让高斯先建立低频结构，再逐步吸收高频纹理。



% METHOD: 2403.11460
\paragraph{Fed3DGS（大规模、城市级与在线重建）（2024.03）}
Fed3DGS~\cite{2403_11460} 关注的问题是：处理大规模、室外或在线 3DGS 重建 Fed3DGS 是宽召回补全中发现的论文，主要根据标题、摘要和 arXiv 类别判定其与 3DGS 的关系。
相对原始 3DGS，Fed3DGS 的变化不是只替换损失函数，而是把表示、初始化、约束或网络预测路径重新组织到任务流程中：该工作被召回是因为论文文本包含 Gaussian/3D Gaussian 相关表述；当前分类先确定其任务域和与 3DGS 的相关性，后续证据流水线再补代码、企业和引用信息。
\textbf{实现细节：} 具体实现可以按输入、核心模块和渲染/优化路径来读。输入端首先取得任务所需的图像、相机、文本、时间序列、语义或几何线索，并把这些线索转换成可以约束高斯中心、尺度、旋转、不透明度、颜色或附加特征的条件。核心结构中，方法使用的主要机制是：从摘要层面看，Fed3DGS 围绕“处理大规模、室外或在线 3DGS 重建”展开。若其正文或代码后续可读，应进一步核对它具体改动的是表示、初始化、渲染、优化、损失、动态建模、压缩或应用接口中的哪一环。 训练或推理时，这些模块通常先给出初始化或中间特征，再通过可微 Gaussian rasterization、重建损失、感知/语义约束、几何正则或任务损失回传到高斯参数；如果方法包含前馈网络，则网络输出的是高斯属性、变形场、压缩码、语义特征或控制变量，而不是只输出最终 2D 图像。渲染路径仍以显式高斯投影、深度排序和 alpha 合成为基础，只是在采样、属性预测、局部编辑、动态变形、压缩编码或下游任务接口处加入额外约束。这样写入本章时，它与前后方法的关系是：相对原始 3DGS 和本章已有锚点方法，它可以理解为 3DGS extends\_or\_uses 宽召回分类阶段的关系判断，后续需要按论文方法部分复核。
\textbf{证据与边界：} 本轮召回分类尚未逐项核对实验表格；后续证据补全阶段需要读取论文源码/PDF或项目页确认数据集、指标和对比方法。 该条来自宽召回批次；源码下载受 arXiv e-print 网络限制，部分论文暂时只能依据标题、摘要和类别分类，需要在入可信书前做人工/LLM精读复核。 本书只采用本地 paper pack、数据库分类、源码/引用登记和论文摘要中可核验的信息；没有本地证据支撑的速度、排行榜或开源完整性不在正文中强化。\textbf{方法信息：} 代码：\url{https://github.com/densoitlab/fed3dgs}；代码关注：densoitlab/fed3dgs has 111 stars (4.0/month, domain-relative cutoff 57.0/month)。
\textbf{设计动机与边界：} 相对原始 3DGS 以独立高斯和图像重建误差为主的设定，Fed3DGS 把问题落在以下可实现的改动上：从摘要层面看，Fed3DGS 围绕“处理大规模、室外或在线 3DGS 重建”展开。若其正文或代码后续可读，应进一步核对它具体改动的是表示、初始化、渲染、优化、损失、动态建模、压缩或应用接口中的哪一环。其设计目的不是单独增加一个模块，而是用上述表示、监督或渲染约束消除原流程在目标任务中的具体失配。该工作被召回是因为论文文本包含 Gaussian/3D Gaussian 相关表述；当前分类先确定其任务域和与 3DGS 的相关性，后续证据流水线再补代码、企业和引用信息。适用时还应保留论文给出的边界：该条来自宽召回批次；源码下载受 arXiv e-print 网络限制，部分论文暂时只能依据标题、摘要和类别分类，需要在入可信书前做人工/LLM精读复核。
% METHOD: 2607.01698
\paragraph{Structure-Aware Gaussian Splatting（结构感知密度与图像频率调度）（2026.07）}
Structure-Aware Gaussian Splatting for Large-Scale Scene Reconstruction~\cite{2607_01698} 回到大场景训练的一个基础矛盾：稀疏初始点往往只含低频结构，而训练图像可能提供高频纹理；若过早用高分辨率图像和强 densification 监督这些区域，3DGS 会用不受控的高斯增殖去拟合缺失结构，产生冗余和伪影。

方法把这个矛盾表述为频率错配：scene frequency bandwidth 由当前高斯尺度和几何支持决定，image sampling frequency 由训练图像分辨率决定。如果后者升得太快，低频或稀疏区域会被高频监督催生出大量不可靠高斯。因此它设计 frequency-aligned resolution and densification scheduler：训练先在低分辨率建立粗全局模型，再根据高斯尺度频率收敛情况逐步提高图像分辨率，并把 densification 的触发与当前结构频率对齐，而不是固定迭代数硬切换。

代码 README 中的流程与论文思路一致：先训练低分辨率 coarse global Gaussian model，再用粗模型做空间分块；每个块从低分辨率开始训练，并根据 SIG dynamic-resolution scheduler 自适应升分辨率；最后合并块高斯并渲染验证视角。方法还引入 Sphere-Constrained Gaussians，把高斯增长约束在初始几何附近，缓解稀疏 COLMAP 点附近的 floaters。相对 LoBE-GS 的负载调度，它调度的是监督频率和密度增长；相对 Mip-Splatting 的抗锯齿，它处理的是大规模稀疏初始化下的训练阶段错配。由于这篇论文很新，代码已经公开但星标和引用还需要后续持续观察。
\textbf{实现细节：} Structure-Aware Gaussian Splatting for Large-Scale Scene Reconstruction。Structure-Aware Gaussian Splatting for Large-Scale Scene Reconstruction（2026.07） 延续 3DGS 在"Structure-aware scheduling for large-scale Gaussian reconstruction"上的基本设定，主要改变是：结构感知调度与频率调制，缓解稀疏初始点导致的失控 densification。3DGS 的显式高斯表示和可微渲染能力保持为基础。其设计动机是：把 densification target signal frequency 和 image sampling frequency 两类调度统一考虑，以控制高频监督和稀疏点初始化之间的矛盾。具体而言，在大规模场景中识别 sparsely observed regions，通过结构感知的 densification 和分辨率/采样调度避免用高频图像监督低频稀疏点时产生 uncontrolled densification。适用边界是：调度策略可能依赖场景结构估计和超参数；极端稀疏区域仍可能需要外部先验。
\textbf{技术实现补充：} Structure-Aware Gaussian Splatting for Large-Scale Scene Reconstruction。Structure-Aware Gaussian Splatting for Large-Scale Scene Reconstruction（2026.07） 延续 3DGS 在"Structure-aware scheduling for large-scale Gaussian reconstruction"上的基本设定，主要改变是：结构感知调度与频率调制，缓解稀疏初始点导致的失控 densification。3DGS 的显式高斯表示和可微渲染能力保持为基础。其设计动机是：把 densification target signal frequency 和 image sampling frequency 两类调度统一考虑，以控制高频监督和稀疏点初始化之间的矛盾。具体而言，在大规模场景中识别 sparsely observed regions，通过结构感知的 densification 和分辨率/采样调度避免用高频图像监督低频稀疏点时产生 uncontrolled densification。适用边界是：调度策略可能依赖场景结构估计和超参数；极端稀疏区域仍可能需要外部先验。
\textbf{方法信息：} 代码：\url{https://github.com/weiyixue999/Signal_Structure_Aware_Gaussian}。

\begin{takeaway}
本章的方法可以分成三条主线。第一条是场景组织：BlockGaussian、Horizon-GS、DroneSplat、OccluGaussian 和 CityGaussian 系列都在回答大场景应该如何分块、分相机、处理遮挡和合并。第二条是表示选择：Octree-GS、Hierarchical 3DGS、PyGS、LODGE、FilterGS 和 NVGS 通过 LoD、过滤或可见性函数让远近尺度不再使用同一套高斯。第三条是系统扩展：DOGS、TideGS、GS-Scale、Grendel、BlitzGS 和多 GPU 抽象把显存、通信和训练并行作为核心问题。在线方法 LongSplat、CL-Splats 与 StreamGS 则把“训练完再渲染”改成“边观察边维护高斯世界”。因此，大规模 3DGS 的技术演进不是单一路径，而是表示层次、训练分治、内存系统和在线更新四个环节一起扩展。
\end{takeaway}
\textbf{设计动机与边界：} 相对原始 3DGS 以独立高斯和图像重建误差为主的设定，Structure-Aware Gaussian Splatting for Large-Scale Scene Reconstruction 把问题落在以下可实现的改动上：在大规模场景中识别 sparsely observed regions，通过结构感知的 densification 和分辨率/采样调度避免用高频图像监督低频稀疏点时产生 uncontrolled densification。其设计目的不是单独增加一个模块，而是用上述表示、监督或渲染约束消除原流程在目标任务中的具体失配。把 densification target signal frequency 和 image sampling frequency 两类调度统一考虑，以控制高频监督和稀疏点初始化之间的矛盾。适用时还应保留论文给出的边界：调度策略可能依赖场景结构估计和超参数；极端稀疏区域仍可能需要外部先验。
